\chapter{BV-BRST formalism and Gaiotto's perspective}
\label{ch:bv-brst}

\label{rem:programme-bv}
\index{Batalin--Vilkovisky|see{BV algebra}}

\section{BV formalism for chiral algebras}

Let $\cA$ be a chiral algebra on a smooth curve~$X$.

\begin{definition}[BV data for chiral algebra]
\index{antibracket}
The BV formalism associates to~$\cA$ the following data:
\begin{enumerate}
\item \emph{Fields.} $\phi \in \mathcal{A}$ (fields of the theory)
\item \emph{Antifields.} $\phi^+ \in \mathcal{A}^*[1]$ (dual shifted by 1)
\item \emph{BV bracket.} $\{\cdot, \cdot\}$ of degree $+1$ (odd Poisson structure)
\item \emph{Action.} $S[\phi, \phi^+]$ satisfying classical master equation
$\{S, S\} = 0$
\end{enumerate}
\end{definition}

\begin{theorem}[BV complex = geometric bar complex; \ClaimStatusProvedElsewhere]
\label{thm:bv-bar-geometric}
The BV complex $(C_{\text{BV}}(\mathcal{A}), Q_{\text{BV}})$ is isomorphic to
the geometric bar complex:
\[C_{\text{BV}}(\mathcal{A}) \cong \bar{B}^{\text{ch}}(\mathcal{A})\]

The BV differential $Q_{\text{BV}} = \{S, -\}$ corresponds to the bar differential.
The precise statement and proof appear in the Costello--Gwilliam framework \cite{CG17}; we give the geometric construction below.
\end{theorem}

\begin{proof}[Geometric construction]
\emph{Step~1: Bar generators.}

The degree-$n$ bar complex is:
\[\bar{B}^n(\mathcal{A}) = \Omega^*(\overline{C}_{n+1}(X), \mathcal{A}^{\boxtimes (n+1)})\]

The generators are operator insertions $\phi_i \in \mathcal{A}$ at marked points, paired with logarithmic $1$-forms $\eta_{ij} = d\log(z_i - z_j)$ dual to the collision divisors.

\emph{Step~2: BV Bracket.}

The BV bracket is a residue operation on the log-de~Rham complex
of~$\overline{C}_{n+1}(X)$.  In local coordinates near the divisor
$D_{ij} = \{z_i = z_j\}$, with $\epsilon = z_i - z_j$:
\[\{\phi(z_i), \eta_{jk}\}
= \delta_{ij}\,\Res_{\epsilon = 0}
\Bigl[\frac{\partial\phi}{\partial z_i}\,
\frac{d\epsilon}{\epsilon(z_i - z_k)}\Bigr]
+ \delta_{ik}\,\Res_{\epsilon = 0}
\Bigl[\frac{\partial\phi}{\partial z_i}\,
\frac{d\epsilon}{\epsilon(z_i - z_j)}\Bigr].\]
Both terms are Poincar\'e residues of logarithmic forms along
normal-crossing divisors of the smooth
variety~$\overline{C}_{n+1}(X)$: the residue map
$\Res_{D_{ij}}\colon \Omega^k_{\log}(\overline{C}_{n+1})
\to \Omega^{k-1}_{D_{ij}}$ is an algebraic morphism of sheaves
on a smooth variety.  The factor $1/(z_i - z_k)$ is a pole along
the transverse divisor $D_{ik}$ (Theorem~\ref{thm:FM}(3):
$D_{ij} \cap D_{ik}$ has normal crossings), so on~$D_{ij}$
it restricts to a regular meromorphic function with poles only
along~$D_{ij} \cap D_{ik}$ --- a codimension-$2$ stratum.
The FM compactification resolves all coincidence singularities
into normal-crossing divisors, and the log-de~Rham complex on
this resolution is the correct algebraic framework: it replaces
distributional regularization with algebraic residues.

The logarithmic $1$-forms $\eta_{ij} = d\log(z_i - z_j)$ generate
$H^1(C_n(X); \mathbb{C})$ by the Arnold--Brieskorn theorem;
for $X = \mathbb{C}$, all $\binom{n}{2}$ classes are linearly independent in degree~$1$
(the Arnold relations constrain products in degree~$\geq 2$), giving
$\dim H^1(C_n(\mathbb{C})) = \binom{n}{2}$.  (For a curve of genus~$g > 0$, additional generators from $H^1(X)$ contribute $2gn$ classes.)

\emph{Step~3: Master Equation.}

The classical master equation $\{S, S\} = 0$ is equivalent to $\dzero^2 = 0$ for the
genus-$0$ bar differential:
\[d = d_{\text{strat}} + d_{\text{int}} + d_{\text{res}}\]
\end{proof}

\subsection{Quantum master equation}

\begin{definition}[BV Laplacian]
The BV Laplacian $\Delta_{\text{BV}}$ is the second-order operator:
\[\Delta_{\text{BV}} = \sum_i \frac{\partial^2}{\partial \phi_i \partial \phi_i^+}\]

In the geometric realization:
\[\Delta_{\text{BV}} = \sum_{i<j} \int \delta(z_i - z_j) 
\frac{\partial}{\partial \eta_{ij}}\]

This inserts delta functions along diagonals—exactly the cobar differential.
\end{definition}

\begin{theorem}[Quantum master equation as Maurer--Cartan equation; \ClaimStatusProvedHere]
\label{thm:qme-bar-cobar}
\index{quantum master equation|textbf}
The quantum master equation
\[\Delta_{\text{BV}} e^{S/\hbar} = 0\]
is equivalent to the Maurer--Cartan equation in the dg Lie algebra
$(\bar{B}^{\ge 1}(\cA),\, \delta,\, [-,-])$: it encodes the
compatibility of bar and cobar differentials under Verdier duality.
This equivalence is natural: it holds functorially for morphisms of
chiral algebras via the BV functor of Theorem~\ref{thm:bv-functor}.
\end{theorem}

\begin{proof}
\emph{Stage~1 (Algebraic equivalence).}
Write $S = S_0 + S_{\text{int}}$ with $S_0$ the free (quadratic) action and $S_{\text{int}}$ the interactions.  Then:
\[\Delta_{\text{BV}} e^{S/\hbar} = \frac{1}{\hbar}\left(\Delta_{\text{BV}} S +
\frac{1}{2\hbar}\{S, S\}\right) e^{S/\hbar}\]

Setting this to zero and using $\{S_0, S_0\} = 0$ (classical master equation for the free theory), $\Delta S_0 = 0$ (no quantum anomaly for quadratic actions), and absorbing $\{S_0, S_{\text{int}}\}$ into a total BV differential $\delta = \Delta_{\text{BV}} + \frac{1}{\hbar}\{S_0, -\}$ gives:
\[\hbar\,\delta\, S_{\text{int}} + \frac{1}{2}\{S_{\text{int}}, S_{\text{int}}\} = 0\]

This is precisely the Maurer--Cartan equation in the dg Lie algebra
$(\bar{B}^{\ge 1}(\cA), \delta, [-,-])$.  The identification of each
term with a bar/cobar operation is provided by Stage~2.

\emph{Stage~2 (Identification of terms via the BV functor).}
The BV functor $\mathrm{BV}: \mathrm{ChirAlg}_X \to \mathrm{BV\text{-}Alg}$
of Theorem~\ref{thm:bv-functor} assigns to each chiral algebra~$\cA$ the BV algebra
$\mathrm{BV}(\cA) = (\bar{B}^{\mathrm{ch}}(\cA), \cdot, \Delta, \{-,-\})$.
Its Verdier duality compatibility (Theorem~\ref{thm:bv-functor}, last part)
establishes that $\mathbb{D}(\bar{B}(\cA)) \cong \Omega(\cA^!)$ as BV algebras,
naturally in~$\cA$.

Under this identification, each BV operation matches a concrete
bar-cobar operation:
\begin{itemize}
\item The bar differential $d_{\bar{B}} = \{S, -\}$ on $\bar{B}(\cA)$ becomes
the cobar differential $d_\Omega$ on $\Omega(\cA^!)$.
\item The BV Laplacian $\Delta$ on $\bar{B}(\cA)$ becomes the cobar product on $\Omega(\cA^!)$.
\item The QME $\hbar\Delta S + \tfrac{1}{2}\{S,S\} = 0$ is the condition ensuring
that both differentials square to zero \emph{simultaneously}, i.e.,
that the BV algebra $(d_{\bar{B}}, \Delta)$ is consistent.
\end{itemize}

Now let $f\colon \cA \to \cB$ be a morphism of chiral algebras.
Functoriality gives a BV morphism
$\mathrm{BV}(f)\colon \mathrm{BV}(\cA) \to \mathrm{BV}(\cB)$
preserving the product, Laplacian, and bracket, while Verdier
duality transports this morphism to the corresponding cobar map on
duals.  Therefore the identification between the quantum master
equation and the bar-cobar compatibility is functorial in~$\cA$; it is
a natural equivalence, not merely a pointwise comparison.
\end{proof}

\begin{remark}[Scope]
Stage~1 uses standard BV algebra manipulations (\ClaimStatusProvedElsewhere,
see Costello--Gwilliam~\cite{CG17}).  Stage~2 assembles the BV
functor (Theorem~\ref{thm:bv-functor}, \ClaimStatusProvedHere) with
the Verdier duality isomorphism $\mathbb{D}(\bar{B}(\cA)) \cong \Omega(\cA^!)$
(Theorem~\ref{thm:verdier-bar-cobar}, \ClaimStatusProvedHere).  The
identification of the QME with the functorial expression of bar-cobar
Verdier compatibility is Theorem~\ref{thm:qme-bar-cobar} above.
\end{remark}

\subsection{The BV--bar dictionary}
\label{subsec:bv-bar-identification}
\index{BV algebra!bar complex identification}

The dictionary is:
\begin{center}
\begin{tabular}{lcl}
\textbf{BV side} && \textbf{Bar complex side} \\[2pt]
fields $\phi \in \cA$ && bar generators $s^{-1}\bar{\cA}$ \\
antibracket $\{-,-\}$ && bracket from simple-pole OPE residue \\
QME: $\hbar\Delta S + \tfrac{1}{2}\{S,S\} = 0$
  && curved $A_\infty$ relation: $m_1^2 = [m_0, -]$ \\
factor $\tfrac{1}{2}$ in QME && Sugawara normalization $m_0 = \kappa/2$
\end{tabular}
\end{center}
The first three entries are reformulations of
Theorem~\ref{thm:config-space-bv} and the proof of
Theorem~\ref{thm:qme-bar-cobar}.  The fourth is the observation
that the QME coefficient $\tfrac{1}{2}$ matches the Sugawara
normalization: the curvature $m_0 = \kappa \cdot \mathbf{1}$ arises
from the double-pole OPE residue, and the factor $\kappa = c/2$ for
the Virasoro algebra (Computation~\ref{comp:virasoro-curvature})
is precisely the $\tfrac{1}{2}$ in the QME applied to the central
charge.

\begin{theorem}[Genus-\texorpdfstring{$0$}{0} amplitudes from bar complex;
\ClaimStatusProvedHere]\label{thm:genus0-amplitude-bar}
\index{string amplitude!bar complex}
Let $\cA$ be a chiral algebra on $\mathbb{P}^1$, and let
$\omega_1, \ldots, \omega_n$ be cocycles in
$\bar{B}^1(\cA)$.  Then the genus-$0$ string amplitude
\[
\mathcal{A}_0(\omega_1, \ldots, \omega_n)
\;=\;
\int_{\overline{\mathcal{M}}_{0,n}}
  \omega_1 \wedge \cdots \wedge \omega_n
\]
equals the bar complex Euler characteristic
$\chi(\bar{B}^n_0(\cA))$, where $\bar{B}^n_0$ denotes the
genus-$0$ component of the degree-$n$ bar complex.
\end{theorem}

\begin{proof}
Both sides compute the same configuration space integral.
The bar complex is realized on
$\overline{C}_n(\mathbb{P}^1) = \overline{\mathcal{M}}_{0,n+1}$
(Theorem~\ref{thm:bar-cobar-isomorphism-main}),
with the bar differential given by residues of the FM propagator
$\omega_{ij} = d\log(z_i - z_j)$ along collision divisors.
The string amplitude $\mathcal{A}_0$ inserts vertex operators
at marked points and integrates the resulting form over the
compactification; this is precisely the definition of the bar
complex pairing.  The Euler characteristic
$\chi(\bar{B}^n_0) = \sum_k (-1)^k \dim H^k(\bar{B}^n_0)$
equals the integral by the Gauss--Bonnet--Chern theorem on
$\overline{\mathcal{M}}_{0,n}$ applied to the logarithmic
de~Rham complex: the logarithmic forms $\eta_{ij}$ provide
a resolution of the constant sheaf on
$\overline{\mathcal{M}}_{0,n}$, and integration of the
top-degree component is the Euler characteristic.
\end{proof}

\begin{remark}[Genus-$1$ partition function]\label{rem:genus1-bv}
At genus~$1$, Theorem~\ref{thm:genus-universality} gives $F_1 = \kappa/24$.  For $\cA = \mathrm{Vir}_c$, $F_1 = c/24 = \operatorname{tr}(\Theta_{\mathrm{Vir}_c}^{(1)}) = (c/2)\cdot(1/12)$, encoding the Mumford isomorphism.
\end{remark}

\begin{remark}[Anomaly = curvature]\label{rem:anomaly-curvature-bv}
\index{anomaly!curvature identification}
By Theorem~\ref{thm:anomaly-koszul}, $d_{\mathrm{bar}}^2 = 0$ for $\cA_{\mathrm{tot}} = \cA_{\mathrm{matter}} \otimes \cA_{\mathrm{ghost}}$ if and only if $\kappa_{\mathrm{tot}} = 0$ (equivalently, $c = 26$ for the bosonic string).  When $\kappa_{\mathrm{tot}} \neq 0$, the universal MC class $\Theta_\cA = \kappa \cdot \eta \otimes \Lambda$ (Theorem~\ref{thm:explicit-theta}) curves the bar complex at every genus.
\end{remark}

\section{Gauge fixing and BRST}

\subsection{BRST from BV}

\begin{definition}[BRST operator]
\index{BRST cohomology|textbf}
The BRST operator $Q_{\text{BRST}}$ arises from gauge fixing the BV action. 
Choose a Lagrangian submanifold $\mathcal{L} \subset$ (fields + antifields):
\[Q_{\text{BRST}} = Q_{\text{BV}}|_{\mathcal{L}}\]

In the chiral algebra context:
\[Q_{\text{BRST}} = Q_0 + Q_1 + Q_2 + \cdots\]
where the decomposition is by \emph{antighost number} (not ghost number): $Q_k$ has antighost number~$k$ and operator dimension~$k-1$.  The total BRST charge has definite ghost number~$+1$.
\end{definition}

\begin{theorem}[BRST cohomology = physical states; \ClaimStatusProvedElsewhere{} \cite{CG17}]
\label{thm:brst-physical-states}
The cohomology of $Q_{\text{BRST}}$ computes physical on-shell states:
\[H^*(Q_{\text{BRST}}) \cong \mathcal{A}_{\text{phys}}\]
\end{theorem}

\begin{example}[Free \texorpdfstring{$bc$}{bc} ghost system]
The $bc$ system has fields $b(z)$ of weight $\lambda$ and $c(z)$ of weight $1-\lambda$, OPE $b(z)c(w)\sim(z-w)^{-1}$, and central charge $c_{bc}=1-3(2\lambda-1)^2$.
For $\lambda=2$, the Virasoro charge is the residue~\eqref{eq:vir-charge-residue} on the domain of Lemma~\ref{lem:vir-domain}, with its specified ghost vacuum and normal ordering.

For the bosonic string, $\lambda = 2$ (so $b$ has weight~$2$ and $c$ has weight~$-1$), giving $c_{\text{ghost}} = 1 - 3(3)^2 = -26$ and ghost stress tensor $T_{\text{ghost}} = -2b\,\partial c - (\partial b)\,c$.

The BRST nilpotence condition:
\[Q^2 = 0 \iff c_{\text{matter}} + c_{\text{ghost}} = 0, \quad \text{i.e., } c_{\text{matter}} = 26\]

For a collision-state differential $d_{\mathrm{res}}$, a comparison map $\Phi$ must satisfy
\[d_{\mathrm{res}}\Phi=\Phi Q_{\mathrm{BRST}}.\]
Theorem~\ref{thm:brst-bar-genus0} constructs such a map only after the target, relative carrier, compatible filtrations, actual symbol quasi-isomorphism, and corrections are supplied.
The augmentation-kernel target for this matter-ghost algebra is excluded by Proposition~\ref{prop:vir-no-augmentation}.
\end{example}

\subsection{Coupling to topological gravity}

\begin{theorem}[Ghost transformation law for log forms;
\ClaimStatusProvedHere]
\label{thm:log-form-ghost-law}
\index{ghost!transformation law}
\index{logarithmic form!ghost transformation}
The logarithmic forms $\eta_{ij} = d\log(z_i - z_j)$
on the configuration space $C_n(X)$ transform under
coordinate changes $z \to w(z)$ exactly as BRST ghosts
for diffeomorphisms:
\[
\eta_{ij} \;\mapsto\; \eta_{ij}
+ d\log\!\left(\frac{w(z_i) - w(z_j)}{z_i - z_j}\right).
\]
Near the diagonal $z_i \to z_j$, the correction term
approaches $d\log w'(z_i)$, which is the standard ghost
cocycle for $\mathrm{Diff}(X)$.  Moreover, the Arnold
relations
$\eta_{ij} \wedge \eta_{jk} + \eta_{jk} \wedge \eta_{ki}
+ \eta_{ki} \wedge \eta_{ij} = 0$
in the Orlik--Solomon algebra provide the quadratic
structure shared by the BRST ghost algebra
($Q_{\mathrm{BRST}}(c) = c\,\partial c$).
\end{theorem}

\begin{proof}
Direct computation: $\eta_{ij} = d\log(z_i - z_j)$
maps to $d\log(w(z_i) - w(z_j))$ under $z \to w(z)$.
Factoring:
\[
d\log(w(z_i) - w(z_j)) = d\log(z_i - z_j)
+ d\log\!\left(\frac{w(z_i) - w(z_j)}{z_i - z_j}\right)
\]
and near the diagonal, L'H\^opital gives
$(w(z_i) - w(z_j))/(z_i - z_j) \to w'(z_i)$, so
the correction approaches $d\log w'(z_i)$.  This is
the BRST ghost 1-cocycle for the Lie algebra
$\mathrm{Vect}(X)$: if $c(z)$ is the ghost field for
$v(z)\partial_z \in \mathrm{Vect}(X)$, then the
infinitesimal version
$\delta_v \eta_{ij} = d(v'(z_i)/(z_i - z_j))$ matches
the BRST variation $\delta c = c\,\partial c$ at leading
order.  The Arnold relations are verified in
Theorem~\ref{thm:arnold-relations}.
\end{proof}

\begin{conjecture}[Bar complex \texorpdfstring{$=$}{=} topological BRST;
\ClaimStatusConjectured]
\label{conj:bar-topological-brst}
The ghost transformation law
(Theorem~\ref{thm:log-form-ghost-law}) extends to a dg
algebra isomorphism:
\[
\bar{B}^{\mathrm{ch}}(\mathcal{A}) \cong
C^*_{\mathrm{top}\text{-}\mathrm{BRST}}(\mathcal{A}
\otimes \mathrm{Diff}(X)).
\]
\end{conjecture}

\begin{remark}[Scope and relation to proved results]
\label{rem:topological-brst-scope}
Theorems~\ref{thm:bar-semi-infinite-km} and~\ref{thm:bar-semi-infinite-w} construct filtered comparisons under their specified coefficient, target, symbol-map, and correction hypotheses.
The string criterion at $c=26$ in Theorem~\ref{thm:brst-bar-genus0} also requires an actual target and relative coordinate carrier.
Propositions~\ref{prop:vir-no-augmentation} and~\ref{prop:vir-relative-target-obstruction} give obstructions to the augmentation-kernel target and the standard restricted transitions.
The topological comparison further requires diffeomorphism ghosts and compatible maps on their full state carriers.

\emph{Homotopy template} (Convention~\ref{conv:hms-levels}): Type~VII (physics-dictionary).  Contributing to Conjecture~\ref{conj:master-bv-brst}.
\end{remark}

\section{The genus-\texorpdfstring{$0$}{0} BRST-bar chain map}
\label{sec:brst-bar-chain-map}
\index{BRST!bar complex chain map|textbf}

\subsection{Semi-infinite cohomology and the BRST complex}

\subsubsection{The state space and the first anomaly}

A degree-one operator can define cohomology only when its degree-two square vanishes on the specified state space.
For the Virasoro algebra, this square is a quadratic ghost operator.
Its coefficient comes from a finite contraction in the ghost representation.
The zero modes then determine which restrictions form cochain complexes.

All vector spaces and tensor products below are algebraic over $\mathbb C$.
For homogeneous operators, the bracket is
$[A,B]_{\mathrm s}=AB-(-1)^{|A||B|}BA$.
Braces denote the bracket of two odd operators.
Matter operators have ghost degree zero.

Let $M$ be a Virasoro module with central charge $C\in\mathbb C$:
\begin{equation}\label{eq:vir-matter-relation}
 [L_m^M,L_n^M]=(m-n)L_{m+n}^M+
 \frac{C}{12}(m^3-m)\delta_{m+n,0}\,1.
\end{equation}
Assume that $M$ is restricted: for each $v\in M$, $L_n^Mv=0$ for all sufficiently large positive $n$.
No inner product or finite-dimensional weight-space hypothesis is imposed.

The ghost Clifford algebra has odd generators $b_n,c_n$, $n\in\mathbb Z$, and relations
\[
 \{b_m,c_n\}=\delta_{m+n,0},\qquad
 \{b_m,b_n\}=\{c_m,c_n\}=0.
\]
Its conformal vacuum $\Omega$ obeys
\begin{equation}\label{eq:vir-ghost-vacuum}
 b_n\Omega=0\quad(n\geq-1),\qquad
 c_n\Omega=0\quad(n\geq2).
\end{equation}
The ghost module is
\[
 \mathcal F=\Lambda\bigl(b_n\ (n\leq-2),\ c_n\ (n\leq1)\bigr)\Omega.
\]
A displayed generator acts by exterior multiplication.
Its Clifford partner acts by left exterior differentiation.
Thus the formulas specify the signs on every monomial.
Set $\deg\Omega=0$, $\deg c_n=1$, and $\deg b_n=-1$.
Set $E\Omega=0$ and assign both $b_n$ and $c_n$ energy $-n$.
The energy of a monomial is the sum of these assignments.

Normal ordering, denoted by colons, moves the creators in~\eqref{eq:vir-ghost-vacuum} to the left.
Each interchange of odd factors contributes a minus sign.
Contraction terms are omitted.
This convention applies to the whole monomial inside the colons.
Define
\begin{align}
 L_m^{\mathrm{gh}}&=\sum_{s\in\mathbb Z}(m+s):b_{m-s}c_s:,
 \label{eq:vir-ghost-L}\\
 R&=-\frac12\sum_{m,n\in\mathbb Z}(m-n):c_{-m}c_{-n}b_{m+n}:,
 \label{eq:vir-ghost-R}\\
 Q&=\sum_{n\in\mathbb Z}c_{-n}L_n^M+R
 \quad\hbox{on }\mathcal D=M\otimes\mathcal F.
 \label{eq:brst-differential}
\end{align}
Tensor factors commute because the matter action is even.
The degree of $Q$ is one.

\begin{lemma}[Algebraic domain]\label{lem:vir-domain}
Every sum in~\eqref{eq:vir-ghost-L}--\eqref{eq:brst-differential} has finitely many nonzero terms on each element of $\mathcal D$.
The same holds for
\[
 \nu_r=\frac12\sum_{a+b=r}(a-b)c_ac_b,
 \qquad
 K=\frac{C-26}{12}\sum_{n>0}(n^3-n)c_{-n}c_n.
\]
These operators and all their finite compositions preserve $\mathcal D$.
\end{lemma}

\begin{proof}
It suffices to use one matter vector and one ghost monomial.
For the matter term, indices $n\leq-2$ require the occupied creator $b_n$.
Only finitely many such indices occur.
The restrictedness of $M$ bounds the remaining positive indices.

For a normally ordered ghost monomial, every annihilator must remove an occupied creator.
Its indices therefore lie in a fixed finite set.
All creator indices have an upper bound: $-2$ for $b$ and $1$ for $c$.
The sum of the indices is fixed in each displayed operator.
Once the annihilator indices are fixed, this sum and the upper bounds give lower bounds for every creator index.
There are at most three factors.
Hence only finitely many terms act nontrivially.
The argument includes terms containing only creators.
It also applies to $\nu_r$ and $K$.
Each output is a finite linear combination of algebraic states, so compositions are defined on the same domain.
\end{proof}

For trivial matter, all $L_n^M$ vanish and $C=0$.
The first nonzero anomalous mode is visible without an infinite summation:
\begin{align*}
 R\Omega&=0, &
 Rb_{-2}\Omega&=-2b_{-2}c_0\Omega-b_{-3}c_1\Omega,\\
 R(b_{-2}c_0\Omega)&=4c_{-2}\Omega+b_{-3}c_0c_1\Omega
       +2b_{-2}c_{-1}c_1\Omega,\\
 R(b_{-3}c_1\Omega)&=5c_{-2}\Omega-2b_{-3}c_0c_1\Omega
       -4b_{-2}c_{-1}c_1\Omega.
\end{align*}
Consequently,
\begin{equation}\label{eq:vir-first-square}
 R^2b_{-2}\Omega=-13c_{-2}\Omega.
\end{equation}
The cancellation of the cubic output leaves a degree-two operator acting on a degree-minus-one input.
The calculation below determines every mode of this operator.

\subsubsection{The ghost central term}

\begin{lemma}[Ghost commutants]\label{lem:vir-commutant}
An operator on $M\otimes\mathcal F$ that supercommutes with every $b_n,c_n$ is zero if its ghost degree is nonzero.
An even operator on $\mathcal F$ commuting with every $b_n,c_n$ is a scalar.
\end{lemma}

\begin{proof}
The common kernel of the vacuum annihilators on $M\otimes\mathcal F$ is $M\otimes\mathbb C\Omega$.
Indeed, these annihilators are all the exterior derivatives of the creator variables.
A finite exterior polynomial killed by all these derivatives is constant.
A supercommuting operator sends this common kernel into itself.
The kernel has ghost degree zero, so an operator of nonzero ghost degree vanishes there.
Supercommutation with the creators makes it vanish on every state.
For an even operator on $\mathcal F$, its value on $\Omega$ is a scalar multiple of $\Omega$.
Commutation with the creators gives the same scalar on all monomials.
\end{proof}

\begin{proposition}[The ghost Virasoro algebra]\label{prop:vir-ghost-central}
On $\mathcal F$,
\begin{align}
 [L_m^{\mathrm{gh}},b_n]&=(m-n)b_{m+n}, &
 [L_m^{\mathrm{gh}},c_n]&=(-2m-n)c_{m+n},
 \label{eq:vir-ghost-transform}\\
 L_0^{\mathrm{gh}}&=E, &
 L_m^{\mathrm{gh}}\Omega&=0\quad(m\geq-1),
 \label{eq:vir-ghost-energy}\\
 [L_m^{\mathrm{gh}},L_n^{\mathrm{gh}}]
 &=(m-n)L_{m+n}^{\mathrm{gh}}
 -\frac{26}{12}(m^3-m)\delta_{m+n,0}\,1.
 \label{eq:vir-ghost-central}
\end{align}
\end{proposition}

\begin{proof}
The elementary contractions are
\[
 [:b_rc_s:,b_n]=\delta_{s+n,0}b_r,
 \qquad
 [:b_rc_s:,c_n]=-\delta_{r+n,0}c_s.
\]
They give~\eqref{eq:vir-ghost-transform}.
A term acting on $\Omega$ needs both factors to be creators.
For $m\geq0$ this is impossible.
For $m=-1$ the only possible term has coefficient zero.
The commutators of $L_0^{\mathrm{gh}}$ with the creators equal those of $E$.
Both operators kill $\Omega$, which proves~\eqref{eq:vir-ghost-energy}.

Let
$A_{m,n}=[L_m^{\mathrm{gh}},L_n^{\mathrm{gh}}]-(m-n)L_{m+n}^{\mathrm{gh}}$.
Substitution of~\eqref{eq:vir-ghost-transform} into the Jacobi identity shows that $A_{m,n}$ commutes with every ghost generator.
For example, its commutator with $b_r$ has coefficient
\[
 (n-r)(m-n-r)-(m-r)(n-m-r)-(m-n)(m+n-r)=0.
\]
For $c_r$, the coefficient is
\[
 (-2n-r)(-2m-n-r)-(-2m-r)(-2n-m-r)
 -(m-n)(-2m-2n-r)=0.
\]
Lemma~\ref{lem:vir-commutant} makes $A_{m,n}$ a scalar.
Its energy is $-m-n$, so this scalar vanishes unless $m+n=0$.

For $m\geq2$, the complete vacuum expression is
\[
 L_{-m}^{\mathrm{gh}}\Omega
 =\sum_{s=2-m}^{1}(s-m)b_{-m-s}c_s\Omega.
\]
Commuting $L_m^{\mathrm{gh}}$ through the first factor gives $(2m+s)b_{-s}$.
This contracts with $c_s$ to give $1$.
Its commutator with the second factor annihilates $\Omega$.
Therefore
\begin{align*}
 A_{m,-m}\Omega
 &=\sum_{s=2-m}^{1}(s-m)(2m+s)\Omega\\
 &=\left(\sum_{s=2-m}^{1}s^2
      +m\sum_{s=2-m}^{1}s-2m^2\cdot m\right)\Omega\\
 &=-\frac{13}{6}(m^3-m)\Omega.
\end{align*}
Here the interval has $m$ integers,
$\sum s=m(3-m)/2$, and
$\sum s^2=m(2m^2-9m+13)/6$.
The cases $m=0,1$ give zero by~\eqref{eq:vir-ghost-energy}.
Antisymmetry gives negative $m$.
This proves~\eqref{eq:vir-ghost-central}.
\end{proof}

Set $\mathcal L_n=L_n^M+L_n^{\mathrm{gh}}$ on $\mathcal D$.
The two summands commute.
Thus~\eqref{eq:vir-matter-relation} and~\eqref{eq:vir-ghost-central} give
\begin{equation}\label{eq:vir-total-central}
 [\mathcal L_m,\mathcal L_n]
 =(m-n)\mathcal L_{m+n}+A(m)\delta_{m+n,0},
 \qquad A(m)=\frac{C-26}{12}(m^3-m).
\end{equation}
This central term determines the square of $Q$ on the full ghost module.

\subsubsection{The square as an operator identity}

\begin{lemma}[Contractions of the charge]\label{lem:vir-Q-contractions}
The operator~\eqref{eq:brst-differential} satisfies
\begin{equation}\label{eq:vir-Q-contractions}
 \{Q,b_n\}=\mathcal L_n,\qquad \{Q,c_n\}=\nu_n.
\end{equation}
Moreover,
\begin{equation}\label{eq:vir-nu-identities}
 [\mathcal L_m,\nu_n]=(-2m-n)\nu_{m+n},\qquad
 [Q,\nu_n]=0.
\end{equation}
\end{lemma}

\begin{proof}
Contracting $b_n$ with the two $c$ factors in $R$ gives
\[
 \{R,b_n\}=\sum_{m\in\mathbb Z}(m-n):c_{-m}b_{m+n}:
 =\sum_{s\in\mathbb Z}(n+s):b_{n-s}c_s:
 =L_n^{\mathrm{gh}}.
\]
Contracting $c_n$ with the $b$ factor gives
\[
 \{R,c_n\}=-\frac12\sum_{r+s=-n}(r-s)c_{-r}c_{-s}=\nu_n.
\]
The matter term contributes $L_n^M$ to the first identity and zero to the second.
All contractions retain the same normal ordering of the uncontracted factors.

The commutator of $\mathcal L_m$ with $\nu_n$ uses~\eqref{eq:vir-ghost-transform}.
After reindexing, the coefficient of $c_uc_v$, $u+v=m+n$, is
\begin{align*}
 \frac12\bigl((u-v-m)(-m-u)+(u-v+m)(-m-v)\bigr)
 &=-\frac12(2m+n)(u-v).
\end{align*}
This proves the first part of~\eqref{eq:vir-nu-identities}.

For the second part, extend $\delta(c_n)=\nu_n$ as an odd derivation on products of $c$ modes.
Their mutual anticommutators are zero.
The coefficient of each alternating cubic monomial in $\delta^2(c_n)$ is a multiple of
\[
 (a-b)(a+b-d)+(b-d)(b+d-a)+(d-a)(d+a-b)=0.
\]
Thus $\delta^2(c_n)=0$.
The second identity in~\eqref{eq:vir-Q-contractions} identifies $[Q,\nu_n]$ with this expression.
Lemma~\ref{lem:vir-domain} permits these expansions on each algebraic input.
\end{proof}

\begin{theorem}[The Virasoro square]\label{thm:vir-square}\label{lem:brst-nilpotence}
For every restricted matter module $M$ as above,
\begin{equation}\label{eq:vir-square}
 \boxed{\displaystyle
 Q^2=\frac{C-26}{12}\sum_{n>0}(n^3-n)c_{-n}c_n
 \quad\text{on }M\otimes\mathcal F.}
\end{equation}
In addition,
\begin{equation}\label{eq:vir-Q-L}
 [Q,\mathcal L_n]=-A(n)c_n,\qquad [Q,\mathcal L_0]=0.
\end{equation}
If $M\neq0$, the operator $Q$ is square-zero on the full state space if and only if $C=26$.
\end{theorem}

\begin{proof}
Put $D_m=[Q,\mathcal L_m]$.
The graded Jacobi identity,~\eqref{eq:vir-Q-contractions}, and~\eqref{eq:vir-total-central} give
\begin{align*}
 \{D_m,b_n\}
 &=[Q,[\mathcal L_m,b_n]]_{\mathrm s}
   -[\mathcal L_m,\{Q,b_n\}]\\
 &=(m-n)\mathcal L_{m+n}-[\mathcal L_m,\mathcal L_n]
 =-A(m)\delta_{m+n,0},\\
 \{D_m,c_n\}
 &=(-2m-n)\nu_{m+n}-[\mathcal L_m,\nu_n]=0.
\end{align*}
Hence $D_m+A(m)c_m$ supercommutes with all ghosts.
It has ghost degree one, so Lemma~\ref{lem:vir-commutant} makes it zero.
This proves~\eqref{eq:vir-Q-L}.

Let $K$ be the right side of~\eqref{eq:vir-square}.
A direct Clifford contraction gives
\[
 [K,b_n]=-A(n)c_n,\qquad [K,c_n]=0.
\]
For instance,
$[c_{-m}c_m,b_n]=\delta_{m+n,0}c_{-m}-\delta_{n-m,0}c_m$.
The oddness of $A$ gives the formula for both signs of $n$.
Also,
\[
 [Q^2,b_n]=[Q,\{Q,b_n\}]=[Q,\mathcal L_n],\qquad
 [Q^2,c_n]=[Q,\nu_n]=0.
\]
Thus $Q^2-K$ commutes with all ghosts and has ghost degree two.
Lemma~\ref{lem:vir-commutant} proves that it is zero.

When $C=26$, the displayed square vanishes.
For nonzero $v\in M$, the converse follows from
\[
 Q^2(v\otimes b_{-2}\Omega)
 =\frac{C-26}{2}\,v\otimes c_{-2}\Omega.
\]
The output is nonzero when $C\neq26$.
\end{proof}

\subsubsection{Fields, zero modes, and the intercept}

Fix a formal coordinate $z$ and define
\[
 b(z)=\sum_n b_nz^{-n-2},\qquad
 c(z)=\sum_n c_nz^{1-n},\qquad
 T^M(z)=\sum_nL_n^Mz^{-n-2}.
\]
The residue extracts the coefficient of $z^{-1}dz$.
For convergent Laurent series it equals the positively oriented integral divided by $2\pi i$.
All uses here are formal.
The two ghost contractions are the expansion of $(z-w)^{-1}$ in nonnegative powers of $w/z$.
For example,
\[
 \sum_{n\geq-1}z^{-n-2}w^{n+1}=\frac1{z-w}.
\]
The ghost stress tensor is
\begin{equation}\label{eq:vir-stress-field}
 T^{\mathrm{gh}}(z)=-2:b(z)\partial c(z):-:\partial b(z)c(z):
 =\sum_mL_m^{\mathrm{gh}}z^{-m-2}.
\end{equation}
Indeed, the coefficient of $b_rc_s$ is $r+2s$.
When $r+s=m$, this is $m+s$, as in~\eqref{eq:vir-ghost-L}.
Equation~\eqref{eq:vir-ghost-transform} specifies the conformal weights $2$ and $-1$.

Multiplying the fields and taking the residue gives exactly
\begin{equation}\label{eq:vir-charge-residue}
 Q=\operatorname{Res}_{z=0}
 \bigl(c(z)T^M(z)+:b(z)c(z)\partial c(z):\bigr)\,dz.
\end{equation}
For the matter term, the residue condition is $m+n=0$.
For the ghost term it is $r+s+t=0$.
Antisymmetrizing the two $c$ factors changes the coefficient $1-t$ to $(s-t)/2$, which gives~\eqref{eq:vir-ghost-R}.
If $\mathcal N$ means full normal ordering of the elementary ghost factors, then
\[
 \tfrac12\mathcal N(cT^{\mathrm{gh}})=:bc\partial c:.
\]
Thus~\eqref{eq:vir-charge-residue} also gives the residue of
$cT^M+\tfrac12\mathcal N(cT^{\mathrm{gh}})$.
Adding $\tfrac32\partial^2c$ to this current leaves its residue unchanged.
This specifies a zero mode, rather than a sum of all current modes.

\begin{proposition}[Intercept and vacuum conventions]\label{prop:vir-intercept}
For $t\in\mathbb C$, set $Q_t=Q+tc_0$.
Then
\begin{align}
 \{Q_t,b_n\}&=\mathcal L_n+t\delta_{n,0},
 & [Q_t,\mathcal L_0]&=0,
 \label{eq:vir-intercept-zero}\\
 Q_t^2&=\sum_{n>0}
 \left(\frac{C-26}{12}(n^3-n)-2tn\right)c_{-n}c_n.
 \label{eq:vir-intercept-square}
\end{align}
For $M\neq0$, absolute nilpotence holds exactly when $C=26$ and $t=0$.

Let $\downarrow=c_1\Omega$ and normal-order relative to this vacuum.
Its annihilators are $b_n$ for $n\geq0$ and $c_n$ for $n\geq1$.
Use a subscript $\downarrow$ for this ordering.
Then
\begin{equation}\label{eq:vir-intercept-dictionary}
 L_{n,\downarrow}^{\mathrm{gh}}=L_n^{\mathrm{gh}}+\delta_{n,0},
 \quad R_{\downarrow}=R+c_0,
 \quad
 \sum_n c_{-n}L_n^M+R_{\downarrow}-ac_0=Q+(1-a)c_0.
\end{equation}
In this convention the nilpotent intercept is $a=1$.
The underlying Clifford module is the same.
The vacuum $\downarrow$ has conformal energy $-1$ and ghost degree one in the $\Omega$ convention.
\end{proposition}

\begin{proof}
The first line follows from~\eqref{eq:vir-Q-contractions} and $[\mathcal L_0,c_0]=0$.
Since $c_0^2=0$,
\[
 Q_t^2=Q^2+t\{Q,c_0\},\qquad
 \{Q,c_0\}=\nu_0=-2\sum_{n>0}n c_{-n}c_n.
\]
This proves~\eqref{eq:vir-intercept-square}.
On $v\otimes\Omega$, its value is $-2t\,v\otimes c_{-1}c_1\Omega$.
Thus absolute nilpotence requires $t=0$.
Theorem~\ref{thm:vir-square} then requires $C=26$.

The change of vacuum reverses only the creator/annihilator pair $(c_1,b_{-1})$.
Its normal products satisfy
\[
 :b_{-1}c_1:_{\downarrow}-:b_{-1}c_1:
 =b_{-1}c_1+c_1b_{-1}=1.
\]
The coefficient of this pair in $L_0^{\mathrm{gh}}$ is $1$, and it occurs in no other $L_n^{\mathrm{gh}}$.
This gives the first identity in~\eqref{eq:vir-intercept-dictionary}.
The same contraction proof as in Lemma~\ref{lem:vir-Q-contractions} gives
\[
 \{R_{\downarrow}-R,b_n\}=\delta_{n,0},\qquad
 \{R_{\downarrow}-R,c_n\}=0.
\]
The difference $R_{\downarrow}-R-c_0$ vanishes by Lemma~\ref{lem:vir-commutant}.
The last identity in~\eqref{eq:vir-intercept-dictionary} follows.
The energy and degree of $\downarrow$ follow from those of $c_1$.
\end{proof}

\begin{proposition}[The nonnegative ghost polarization]\label{prop:vir-polarization-map}
Set $c^m=c_{-m}$, so that $\{b_m,c^n\}=\delta_m^n$.
Let $\mathcal F_{\mathrm{pol}}$ have vacuum $\Omega_{\mathrm{pol}}$ killed by $c^m$ for $m\leq0$ and by $b_m$ for $m>0$.
Its creators are $c^m$ for $m>0$ and $b_m$ for $m\leq0$.
Give its vacuum degree and energy zero, and assign these creators degrees $1,-1$ and energies $m,-m$ respectively.
There is a Clifford-module isomorphism
\[
 I:\mathcal F_{\mathrm{pol}}\longrightarrow\mathcal F,
 \qquad I(x\Omega_{\mathrm{pol}})=x\,c_0c_1\Omega,
\]
where $x$ is a product of the indicated creators.
It raises ghost degree by $2$ and lowers energy by $1$.
Let $R_{\mathrm{pol}}$ denote the cubic expression~\eqref{eq:vir-ghost-R} with this ordering.
Under $1_M\otimes I$,
\[
 \sum_m c^mL_m^M+R_{\mathrm{pol}}-ac^0
 \quad\longmapsto\quad Q+(1-a)c_0.
\]
Thus its critical charge term is $-c^0$, with both degree and energy shifts specified by $I$.
\end{proposition}

\begin{proof}
The vector $\Xi=c_0c_1\Omega$ obeys the required annihilator conditions.
For $m\leq0$, $c_{-m}\Xi=0$: $c_0$ and $c_1$ square to zero, and all $c_n$ with $n\geq2$ annihilate $\Xi$.
For $m>0$, $b_m\Xi=0$ by the Clifford relations and~\eqref{eq:vir-ghost-vacuum}.
Hence the displayed map intertwines the Clifford actions.
Away from the pairs $(b_0,c_0)$ and $(b_{-1},c_1)$ it identifies the same exterior creators.
On each of these two pairs, exterior multiplication and exterior differentiation exchange the empty and occupied states, with nonzero signs.
Their four-state tensor product is invertible.
This proves that $I$ is an isomorphism.
Since $\Xi$ has degree $2$ and energy $-1$, the grading statements follow.

Relative to $\Omega$, this ordering reverses both specified pairs.
The pair $(b_{-1},c_1)$ contributes $1$ to $L_0^{\mathrm{gh}}$, as in Proposition~\ref{prop:vir-intercept}.
The pair $(b_0,c_0)$ has coefficient zero there.
Consequently the transported ghost stress modes are $L_n^{\mathrm{gh}}+\delta_{n,0}$.
The contraction identities and Lemma~\ref{lem:vir-commutant} then give
$I R_{\mathrm{pol}}I^{-1}=R+c_0$.
The matter term commutes with $I$, which proves the charge formula.
\end{proof}

\subsubsection{Absolute and relative complexes}

Here absolute and relative refer to the $b_0$ constraint.
The central character of the Virasoro module is fixed, and the ghost variables are indexed by the modes $L_n$.
There is no ghost for the central generator in this carrier.

Set
\[
 \mathcal D_t^{\mathrm{rel}}=
 \ker(b_0:\mathcal D\to\mathcal D)
 \cap\ker(\mathcal L_0+t:\mathcal D\to\mathcal D).
\]
It inherits the ghost grading of $\mathcal D$.
The first identity in~\eqref{eq:vir-intercept-zero} and the commutation with $\mathcal L_0$ show that $Q_t$ preserves this subspace.
On it, the square is the restriction of~\eqref{eq:vir-intercept-square}.

\begin{corollary}[The critical complexes]\label{cor:vir-complexes}
At $C=26$ and $t=0$, the pairs
\[
 C_{\mathrm{abs}}^q(M)=M\otimes\mathcal F^q,
 \qquad
 C_{\mathrm{rel}}^q(M)=\mathcal D_0^{\mathrm{rel}}\cap(M\otimes\mathcal F^q)
\]
with differential $Q$ are cochain complexes of complex vector spaces.
The inclusion $C_{\mathrm{rel}}^\bullet(M)\to C_{\mathrm{abs}}^\bullet(M)$ is a chain map.
A homomorphism of restricted Virasoro modules induces chain maps by tensoring with $1_{\mathcal F}$.
If $\mathcal L_0$ acts semisimply, every nonzero weight $\lambda$ subcomplex of the absolute complex is contractible by $\lambda^{-1}b_0$.
\end{corollary}

\begin{proof}
The domain lemma defines the differential in every degree.
Theorem~\ref{thm:vir-square} gives its square, and~\eqref{eq:vir-intercept-zero} gives stability of the relative subspace.
A module homomorphism commutes termwise with~\eqref{eq:brst-differential}, hence with the finite sums on each state.
On the weight-$\lambda$ subspace,
$Q(\lambda^{-1}b_0)+(\lambda^{-1}b_0)Q=1$ by~\eqref{eq:vir-Q-contractions}.
\end{proof}

Relative nilpotence depends on the actual relative carrier.
The full-state converse in Theorem~\ref{thm:vir-square} cannot be restricted without checking its detecting vectors.

\begin{proposition}[Nonnegative matter weights]\label{prop:vir-relative-three-term}
Suppose $M=\bigoplus_{j\geq0}M_j$, with $L_0^M|_{M_j}=j$ and $L_n^MM_j\subseteq M_{j-n}$.
At $t=0$, for every central charge $C$, the relative carrier is
\begin{equation}\label{eq:vir-relative-carrier}
 \mathcal D_0^{\mathrm{rel}}
 =(M_0\otimes\mathbb C\Omega)
 \oplus(M_1\otimes\mathbb Cc_1\Omega)
 \oplus(M_0\otimes\mathbb Cc_{-1}c_1\Omega).
\end{equation}
Its differential is the three-term complex
\begin{equation}\label{eq:vir-relative-three-term}
 M_0\xrightarrow{\ L_{-1}^M\ } M_1
 \xrightarrow{\ L_1^M\ }M_0
\end{equation}
in ghost degrees $0,1,2$.
In particular, $Q^2=0$ on this relative carrier for every $C$.
\end{proposition}

\begin{proof}
The only negative ghost energy is that of $c_1$, and it can occur at most once.
Thus $\mathcal F$ has minimum energy $-1$.
Every $b$ creator has energy at least $2$, so no $b$ creator occurs at ghost energy $0$ or $-1$.
The complete bases at these two energies are
\[
 \mathcal F_{-1}=\langle c_1\Omega,c_0c_1\Omega\rangle,
 \qquad
 \mathcal F_0=\langle\Omega,c_0\Omega,c_{-1}c_1\Omega,c_{-1}c_0c_1\Omega\rangle.
\]
The condition $b_0=0$ removes precisely the terms containing $c_0$.
The total weight-zero condition therefore gives~\eqref{eq:vir-relative-carrier}.

The matter term and ghost term of~\eqref{eq:brst-differential} give
\begin{align*}
 Q(v\otimes\Omega)&=L_{-1}^Mv\otimes c_1\Omega &&(v\in M_0),\\
 Q(w\otimes c_1\Omega)&=L_1^Mw\otimes c_{-1}c_1\Omega &&(w\in M_1),\\
 Q(v\otimes c_{-1}c_1\Omega)&=0 &&(v\in M_0).
\end{align*}
In the middle line, $c_0L_0^Mw\otimes c_1\Omega$ cancels $w\otimes Rc_1\Omega=-w\otimes c_0c_1\Omega$.
For the last line, $R(c_{-1}c_1\Omega)=0$ by the two contractions in Lemma~\ref{lem:vir-Q-contractions}.
Finally,
$L_1^ML_{-1}^Mv=2L_0^Mv+L_{-1}^ML_1^Mv=0$ on $M_0$.
This proves~\eqref{eq:vir-relative-three-term} and its nilpotence.
\end{proof}

For example, the universal Virasoro vacuum module of central charge $C$ is induced from a vector $\mathbf1$ killed by $L_n$ for $n\geq-1$.
The central generator acts by $C$.
The ordered monomials in $L_{-n}$, $n\geq2$, form its basis by the Poincar\'e--Birkhoff--Witt theorem.
Hence $M_0=\mathbb C\mathbf1$ and $M_1=0$.
Its relative complex has basis $\mathbf1\otimes\Omega,\mathbf1\otimes c_{-1}c_1\Omega$ and zero differential for every $C$.
Trivial matter gives the same relative complex at $C=0$.
Its absolute square is nonzero by~\eqref{eq:vir-first-square}.
At the other boundary, suppose a restricted module contains $v\neq0$ with $L_0^Mv=-2v$.
Then $v\otimes b_{-2}\Omega$ lies in $\mathcal D_0^{\mathrm{rel}}$ and
\[
 Q^2(v\otimes b_{-2}\Omega)=\frac{C-26}{2}v\otimes c_{-2}\Omega.
\]
This relative square is nonzero when $C\neq26$.
Thus the critical central charge guarantees relative nilpotence for every restricted module, while particular relative carriers can have further vanishing.

\subsubsection{Conformal vertex algebra states}

Let $V$ be an ordinary conformal vertex algebra over $\mathbb C$ with conformal vector $\omega$.
Write
$Y(\omega,z)=\sum_nL_n^Vz^{-n-2}$ and let $C$ be its central charge.
The lower truncation axiom makes $V$ a restricted matter module.
Consequently~\eqref{eq:brst-differential} defines the degree-one operator on the full algebraic state carrier $V\otimes\mathcal F$.
Equation~\eqref{eq:vir-charge-residue} identifies it with the zero mode of the specified normal-ordered current.
At $C=26$, Corollary~\ref{cor:vir-complexes} supplies both the absolute complex and its $b_0$-relative subcomplex.
When $V$ has nonnegative integral conformal weights, Proposition~\ref{prop:vir-relative-three-term} computes the relative complex directly.

The ghost degree here is the cohomological degree.
Conformal weight is the eigenvalue of $L_0^V+L_0^{\mathrm{gh}}$ and is independent of ghost degree.
The residue uses the fixed formal coordinate $z$.
A comparison with a complex of collision-state sheaves requires maps in that sheaf category, with compatible restrictions and differentials.
The state-space construction above supplies its Virasoro differential and its square.

For the conformal vertex algebra $\cA$, set $V=\cA$, $C=c(\cA)$, and $\cA_{\mathrm{tot}}=V\otimes\mathcal F$ on algebraic states.
The notations $Q_{\mathrm{BRST}}=Q$ and $L_0^{\mathrm{tot}}=\mathcal L_0$ use the conformal vacuum and zero intercept.
At $c(\cA)=26$, define
\[
 H^{\infty/2+q}(\cA)=H^q(C_{\mathrm{abs}}^\bullet(\cA),Q).
\]
The $b_0$-relative cohomology uses the separate complex $C_{\mathrm{rel}}^\bullet(\cA)$.

\begin{remark}[The Virasoro anomaly and a bar comparison]
\label{rem:brst-nilpotence-periodicity}
The quadratic ghost operator in Theorem~\ref{thm:vir-square} has coefficient $c(\cA)-26$.
A comparison with a bar differential requires an actual chain map on its specified collision-state carrier.
If anomaly operators are also to be identified, that map must intertwine those operators.
Equality of their scalar parameters alone does not supply such a map.
\end{remark}

\subsection{Augmentations and the relative coordinate carrier}

The critical Virasoro charge acts on an algebraic state space.
A comparison with chiral chains also needs a target and compatible coordinate transitions.
Two finite identities obstruct the augmentation-kernel target and the standard restricted transitions.

\begin{proposition}[The augmentation-kernel obstruction]\label{prop:vir-no-augmentation}
Let $V\neq0$ be a unital vertex algebra over $\mathbb C$, and let $A=V\otimes bc$ use the ghost system above.
There is no unital vertex superalgebra homomorphism $A\to\mathbb C$ with trivial target vertex products.
Consequently the reduced chiral construction whose coefficient sheaf is a chiral augmentation kernel cannot be instantiated on $A$ with its original collision products.
\end{proposition}

\begin{proof}
Write $\Omega_A=\mathbf1_V\otimes\Omega$, $b=\mathbf1_V\otimes b_{-2}\Omega$, and $c=\mathbf1_V\otimes c_1\Omega$ in $A$.
With the vertex convention $Y(a,z)=\sum_r a_{(r)}z^{-r-1}$,
\[
 b_{(0)}c=\mathbf1_V\otimes b_{-1}c_1\Omega=\Omega_A.
\]
Every nonnegative product vanishes in the trivial vertex algebra.
A unital homomorphism $\varepsilon$ would therefore give
$1=\varepsilon(b_{(0)}c)=\varepsilon(b)_{(0)}\varepsilon(c)=0$.
For $V=\operatorname{Vir}_{26}$, set $T=(L_{-2}^V\mathbf1_V)\otimes\Omega$.
There is also the independent matter identity
\[
 T_{(3)}T=(L_2^VL_{-2}^V\mathbf1_V)\otimes\Omega=13\Omega_A.
\]

On a coordinate disk, let $t=z_1-z_2$ and normalize
$\operatorname{Res}(t^{-1}dt)=1$.
The full collision maps include
\[
 \operatorname{Res}_{t=0}Y(b,t)c\,dt=\Omega_A,
 \qquad
 \operatorname{Res}_{t=0}t^3Y(T,t)T\,dt=13\Omega_A.
\]
A chiral augmentation to the unit chiral algebra must preserve these residues, which gives the same contradiction on that disk.
It therefore cannot exist globally.
The forms $dt=t\,d\log t$ and $t^3dt=t^4d\log t$ have regular logarithmic coefficients.
Restricting to such coefficients does not remove the obstruction.

The linear vacuum projection $\pi:A\to\mathbb C$ has $b,c\in\ker\pi$ but $b_{(0)}c\notin\ker\pi$.
Its kernel is not closed under the original collision maps.
The quotient by the vacuum also fails to inherit all vertex products, since
$(\Omega_A)_{(-1)}b=b$ whereas $0_{(-1)}b=0$.
Thus neither linear operation supplies the missing chiral augmentation.
\end{proof}

\begin{proposition}[Relative states, absolute states, and coordinate changes]
\label{prop:vir-relative-target-obstruction}
Take $V=\operatorname{Vir}_{26}$ and the zero-intercept charge of Theorem~\ref{thm:vir-square}.
In this proposition, $\Omega$ denotes the tensor-product vacuum $\mathbf1_V\otimes\Omega_{\mathrm{gh}}$.
Put $H=L_0^V+L_0^{\mathrm{gh}}$ and
\[
 e=\Omega,\quad u=c_0\Omega,\quad
 v=c_{-1}c_1\Omega,\quad w=c_{-1}c_0c_1\Omega.
\]
The weight-zero absolute complex has these four basis states and differential
\[
 Qe=0,\qquad Qu=-2v,\qquad Qv=Qw=0.
\]
The relative complex is $\mathbb Ce\oplus\mathbb Cv$ with zero differential.
The absolute cohomology has one class in degrees zero and three, and no others.
The inclusion of relative states into absolute states kills $[v]$ and is not a quasi-isomorphism.

Fix the convention that the coordinate flow $z\mapsto z/(1+sz)$ acts by
$g_s=\exp(s\mathcal L_1)$, where $\mathcal L_1=L_1^V+L_1^{\mathrm{gh}}$.
Then
\begin{equation}\label{eq:vir-relative-coordinate-obstruction}
 g_sv=v-sc_0c_1\Omega,
 \qquad b_0g_sv=-sc_1\Omega\neq0\quad(s\neq0).
\end{equation}
The fixed relative fiber therefore does not define a subbundle through restriction of these state-bundle transitions.
\end{proposition}

\begin{proof}
The universal vacuum module has $V_0=\mathbb C\mathbf1$ and $V_1=0$.
The energy argument of Proposition~\ref{prop:vir-relative-three-term} gives exactly the four displayed absolute states and the two relative states.
The charge contraction gives $Qu=-2v$.
Nilpotence gives $Qv=0$, and weight zero has no degree-four state, so $Qw=0$.
At every nonzero total weight $\lambda$, the homotopy $b_0/\lambda$ contracts the absolute complex by Corollary~\ref{cor:vir-complexes}.
On the four-state complex, the homotopy $h(v)=-u/2$, zero on $e,u,w$, contracts it onto $\mathbb Ce\oplus\mathbb Cw$.
All state sums are algebraic, so these homotopies define maps on the full carrier.
In particular, $v=Q(-u/2)$ in the absolute complex.

The ghost transformation law~\eqref{eq:vir-ghost-transform} gives
\[
 \mathcal L_1v=-c_0c_1\Omega,
 \qquad \mathcal L_1(c_0c_1\Omega)=0.
\]
The exponential on $v$ consequently terminates after its linear term.
Since $b_0c_0c_1\Omega=c_1\Omega$, this proves~\eqref{eq:vir-relative-coordinate-obstruction}.
Also $Hg_sv=sc_0c_1\Omega$, so the same transition fails the weight-zero condition.
A restricted transition must preserve the model fiber, which fails here.
\end{proof}

For a fixed coordinate, the relative vector-space complex remains defined.
A sheaf comparison requires a separately constructed relative sheaf complex or a family of transported constraints with compatible transition maps.
For source and target sheaf complexes $\mathcal C,\mathcal T$, local maps must satisfy
\[
 d_{\mathcal T}\Phi_U=\Phi_Ud_{\mathcal C},\qquad
 \operatorname{res}^{\mathcal T}_{UW}\Phi_U
 =\Phi_W\operatorname{res}^{\mathcal C}_{UW}\quad(W\subset U),
\]
and $g^{\mathcal T}_{ij}\Phi_j=\Phi_i g^{\mathcal C}_{ij}$ in coordinate trivializations.
Ordinary or derived global sections require these complexes and their transitions first.
The restricted transitions in~\eqref{eq:vir-relative-coordinate-obstruction} do not supply them.

A different unital chiral-chain construction may avoid an augmentation kernel.
Its coefficient complexes, vacuum insertions, all collision maps, differential, and section functor must be constructed before it can serve as a comparison target.
For the relative vacuum complex, a degree-preserving quasi-isomorphism requires target cohomology precisely in degrees zero and two.
Proposition~\ref{prop:vir-relative-target-obstruction} excludes the absolute state complex as a replacement target.
These assertions concern the specified augmentation-kernel construction and state transitions; they do not exclude every chiral-chain construction.

\subsection{The chain map}

\begin{theorem}[Filtered comparison with a relative string BRST complex]
\label{thm:brst-bar-genus0}
\index{BRST!bar complex quasi-isomorphism}
Let $\cA$ be a conformal vertex algebra with $c(\cA)=26$, and fix a coordinate for its algebraic states. Set $\cA_{\mathrm{tot}}=\cA\otimes bc$.
Choose a graded $Q$-stable algebraic domain
$\mathcal D_{\mathrm{rel}}\subseteq\ker b_0\cap\ker L_0^{\mathrm{tot}}$
in the matter and ghost state space of Theorem~\ref{thm:vir-square}.
Use the restriction of the operator~\eqref{eq:brst-differential} with its conformal vacuum and zero intercept.
Theorem~\ref{thm:vir-square} gives $Q^2=0$ on this domain.
Let $T=(T^\bullet,d_T)$ be a separately specified square-zero cochain complex of complex vector spaces.
A chiral-chain application requires its actual coefficient complexes, vacuum normalization, collision maps, and prescribed section functor.
A global relative application also requires the source descent and section-carrier identifications described after Proposition~\ref{prop:vir-relative-target-obstruction}.
The augmentation-kernel target for $\cA_{\mathrm{tot}}$ is excluded by Proposition~\ref{prop:vir-no-augmentation}; the statement makes no such target assignment.
Write $C=(\mathcal D_{\mathrm{rel}},Q)$.
Assume $C,T$ have exhaustive increasing filtrations with $F_{-1}=0$, preserved by their differentials.
Choose degree-preserving splittings by filtration order and an actual graded chain quasi-isomorphism
\begin{equation}\label{eq:arkhipov-classical}
\Phi_0:\operatorname{gr}C\longrightarrow\operatorname{gr}T.
\end{equation}
Suppose the obstruction classes of Theorem~\ref{thm:fl19-lift-obstruction} vanish along compatible choices of corrections in the Hom complexes of these carriers.
Then the sum of those corrections is a filtered homotopy equivalence
\[
\Phi=\sum_{r\ge0}\Phi_r:
(\mathcal D_{\mathrm{rel}},Q)\longrightarrow
(T^\bullet,d_T).
\]
In particular, it induces an isomorphism on cohomology.
The complete version requires the explicitly split products and filtration-raising maps specified after~\eqref{eq:fl19-cone-contraction}.
\end{theorem}
\begin{proof}
Resolve both differentials by their filtration orders:
$Q=\sum_{i\ge0}Q_i$ and $d_T=\sum_{i\ge0}d_{T,i}$.
On maps of degree $a$, the order-zero Hom differential is
$\partial_0u=d_{T,0}u-(-1)^a uQ_0$.
After choosing $\Phi_0,\ldots,\Phi_{r-1}$, form the actual degree-one defect
\[
\omega_r=\sum_{i=1}^{r}
(d_{T,i}\Phi_{r-i}-\Phi_{r-i}Q_i).
\]
The square-zero identities and the lower-order equations give $\partial_0\omega_r=0$.
The assumed vanishing gives a map $\Phi_r$ with $\partial_0\Phi_r=-\omega_r$.
For the chosen splittings of the graded coefficient complexes, Lemma~\ref{lem:fl19-hom-primitive} gives the primitive $-S(\omega_r)$ whenever its cohomology projection vanishes.
The choices must continue to satisfy the higher obstruction equations.

For an input of order $p$, only $\Phi_0,\ldots,\Phi_p$ act nontrivially.
Their component equations imply $d_T\Phi=\Phi Q$ on every input.
The associated graded map is the specified $\Phi_0$.
The cone of $\Phi$ therefore has an acyclic graded complex.
Its explicit contraction~\eqref{eq:fl19-cone-contraction} supplies a filtered inverse and both chain homotopies.
On the stated complete products, the same equations hold componentwise and the sums converge in their specified filtration.
\end{proof}

\begin{remark}[Compatibility with the BV constructions]
\label{rem:bv-convergence}
\index{BV complex!convergence with bar complex}
Theorem~\ref{thm:brst-bar-genus0} constructs a comparison of the specified filtered cochain complexes.
Compatibility with the BV operations in~\eqref{eq:master-square}, Verdier duality, or the genus-one curvature requires the symbol map and every correction to intertwine those operations.
A chosen vector-space contraction does not impose these additional equations.
\end{remark}

\begin{remark}[Low-degree assignments and the relative domain]
\label{rem:brst-bar-explicit}
Any proposed low-degree formula for $\Phi$ must start with a vector in $\mathcal D_{\mathrm{rel}}$ and satisfy $d_T\Phi(v)=\Phi(Qv)$.
The relative condition already excludes some ghost monomials.
For a vacuum $|0\rangle$ with $b_0|0\rangle=0$ and $c_1|0\rangle\ne0$, the ghost relations give
\[
 b_0c_0c_1|0\rangle=c_1|0\rangle\ne0,
\]
since $\{b_0,c_0\}=1$ and $\{b_0,c_1\}=0$.
Thus $c_0c_1|0\rangle$ is outside $\ker b_0$ and cannot be assigned an image by a map defined only on the relative complex.
Assignments of ghost states to logarithmic collision forms require the actual source domain, target coefficient sheaf, and chain equation.
The recursion in the theorem constructs higher terms once those initial data and compatible primitives are supplied.
\end{remark}

\begin{remark}[Filtration and coefficient hypotheses]
\label{rem:brst-bar-scope}
The state-order filtration must include every matter, ghost, form, and collision factor in its actual carrier.
Preservation of conformal energy does not identify their symbols with a zero-mode Lie algebra complex.
Example~\ref{ex:fl19-no-lift} shows that an isomorphism of first pages does not construct a filtered chain lift.
Algebraic direct sums and complete products require the distinct hypotheses stated in Theorem~\ref{thm:brst-bar-genus0}.
The central-charge condition $c(\cA)=26$ alone supplies neither the target nor its symbol map and compatible corrections.
For $\cA\otimes bc$, Proposition~\ref{prop:vir-no-augmentation} excludes the native augmentation-kernel target.
For the universal Virasoro vacuum, Proposition~\ref{prop:vir-relative-target-obstruction} also excludes restriction of the standard coordinate transitions to the relative fiber.
\end{remark}

\begin{corollary}[Virasoro anomaly cancellation and chain comparison;
\ClaimStatusProvedHere]
\label{cor:anomaly-physical-genus0}
\index{anomaly cancellation!BRST proof}
Let $\cA\neq0$ be a conformal vertex algebra of central charge $c$, and use the conformal ghost vacuum.
Then the absolute Virasoro charge has square
\[
 Q_{\mathrm{BRST}}^2=\frac{c-26}{12}\sum_{n>0}(n^3-n)c_{-n}c_n.
\]
Thus $Q_{\mathrm{BRST}}^2=0$ on the full state space if and only if $c=26$.
The condition $c+c_{\mathrm{gh}}=0$ is equivalent to this vanishing, since $c_{\mathrm{gh}}=-26$.
For the scalar normalization $\kappa_{\mathrm{tot}}=(c-26)/2$, the same condition is $\kappa_{\mathrm{tot}}=0$.
At $c=26$ the relative complex is defined by Corollary~\ref{cor:vir-complexes}.
Under all separately specified target, source-descent, section-carrier, filtration, symbol-map, and compatible-correction hypotheses of Theorem~\ref{thm:brst-bar-genus0}, its map $\Phi$ identifies the cohomology of the two specified complexes.
This criterion does not instantiate an augmentation-kernel bar of $\cA\otimes bc$, which Proposition~\ref{prop:vir-no-augmentation} excludes.
Transport of anomaly operators requires the additional equality
$\Phi Q^2=B\Phi$ for the supplied target anomaly operator $B$.
\end{corollary}

\begin{proof}
The square and its converse are Theorem~\ref{thm:vir-square}.
Proposition~\ref{prop:vir-ghost-central} supplies the ghost central charge.
The stated scalar normalization gives the numerical equivalence.
Corollary~\ref{cor:vir-complexes} gives the relative complex at $c=26$.
Theorem~\ref{thm:brst-bar-genus0} constructs the chain quasi-isomorphism under its full hypotheses.
The equality $\Phi Q^2=B\Phi$ is exactly compatibility with the anomaly operators.
For nonnegative matter weights, Proposition~\ref{prop:vir-relative-three-term} shows why relative nilpotence alone cannot imply $c=26$.
\end{proof}

\begin{remark}[Higher genus]
\label{rem:brst-bar-higher-genus}
At genus $g \geq 1$, the BRST complex requires additional data:
the path integral measure on $\overline{\mathcal{M}}_g$ and the
Costello renormalization framework~\cite{costello-renormalization} for handling
ultraviolet divergences.  The bar complex side is well-defined at
all genera via configuration spaces on $\Sigma_g$
(Chapter~\ref{chap:higher-genus}).  The chain map $\Phi$ should
extend to a map of genus-graded complexes, with the genus-$g$
component involving Costello's counterterm expansion.  This
extension is the subject of Programme~VI-a in
\S\ref{sec:modular-koszul-programme}.
\end{remark}

\section{Semi-infinite cohomology and the bar complex}
\label{sec:semi-infinite-bar}
\index{semi-infinite cohomology|textbf}

\subsection{Semi-infinite cohomology}

\begin{definition}[A specified semi-infinite complex]
\label{def:semi-infinite-cohomology}
Let $\mathfrak l=\bigoplus_m\mathfrak l_m$ be a graded Lie algebra with finite-dimensional components and let $M$ be a coefficient module.
Choose a polarization, its ghost space $\mathcal F$, and a locally finite degree-one operator $Q_{\mathrm{si}}$ on a specified graded algebraic domain $\mathcal D\subseteq M\otimes\mathcal F$, satisfying $Q_{\mathrm{si}}^2=0$ on that domain.
For an ordinary module with the nonnegative energy grading of Lemma~\ref{lem:fl20-ghost-domain}, the stated current and normal-ordered cubic terms act on $\mathcal D=M\otimes\mathcal F$.
Their square-zero identity remains a separate requirement.
The relative ghost degree labels $C^{\infty/2+n}=\mathcal D^n$, and
$H^{\infty/2+n}(\mathfrak l,M)=H^n(\mathcal D,Q_{\mathrm{si}})$.
The treatment of central generators and any charge term are part of these data.
The creation and annihilation conventions, and their significance, are given in Example~\ref{ex:fl19-central-ghost} and the construction preceding it.
\end{definition}

\begin{remark}[Semi-infinite and string differentials]
\label{rem:semi-infinite-vs-brst}
A ghost polarization does not by itself cancel an anomaly.
Example~\ref{ex:fl19-central-ghost} gives a nonzero square when a central term is retained in the current action but omitted from the ghost differential.
Including the central ghost gives a different, acyclic absolute complex at nonzero level in that example.
An identification with a string BRST complex therefore requires the actual coefficient action, central treatment and differential, in addition to a comparison of graded vector spaces.
\end{remark}

\subsection{The bar--semi-infinite identification for Kac--Moody algebras}

\begin{theorem}[Filtered comparison with a specified semi-infinite target]
\label{thm:bar-semi-infinite-km}
Let $C=\bar B^{\mathrm{ch}}(\widehat{\fg}_k)$ be a supplied native cochain complex on full collision-state sheaves and their specified section carrier.
Choose the graded current Lie algebra $\mathfrak l$, its coefficient action on $V_k$, its ghost polarization and central treatment. Let $T=C^{\infty/2+\bullet}(\mathfrak l,V_k)$ denote the resulting specified square-zero complex of Definition~\ref{def:semi-infinite-cohomology}.
Suppose their increasing filtrations start at order zero and are exhaustive, and that their differentials preserve them.
Suppose an actual graded chain quasi-isomorphism $f_0:\operatorname{gr}C\to\operatorname{gr}T$ is given.
If the obstruction classes of Theorem~\ref{thm:fl19-lift-obstruction} vanish along compatible choices of corrections, then the resulting sum defines a filtered homotopy equivalence
\begin{equation}\label{eq:bar-semi-infinite}
\Psi=\sum_{r\ge0}f_r:C\longrightarrow T.
\end{equation}
In particular,
\begin{equation}\label{eq:bar-equals-semi-infinite}
H^*(\bar B^{\mathrm{ch}}(\widehat{\fg}_k))
\cong H^{\infty/2+*}(\mathfrak l,V_k).
\end{equation}
The same conclusion holds on the explicitly split complete carriers described after~\eqref{eq:fl19-cone-contraction}.
\end{theorem}
\begin{proof}
At order $r$, calculate the defect
$\omega_r=\sum_{i=1}^r(d_{T,i}f_{r-i}-f_{r-i}d_{C,i})$ on the actual coefficient complexes.
The hypothesis supplies $f_r$ with $\partial_0 f_r=-\omega_r$.
When its cohomology projection vanishes, Lemma~\ref{lem:fl19-hom-primitive} gives the explicit choice $f_r=-S(\omega_r)$.
The sum in~\eqref{eq:bar-semi-infinite} is finite on every input of finite order, and the component equations prove its chain identity.
Its graded map is $f_0$.
The cone contraction~\eqref{eq:fl19-cone-contraction} gives the inverse and both homotopies, including their signs.
The complete case uses componentwise convergent sums on the specified products.
\end{proof}

\begin{remark}[The coefficient comparison required by the theorem]
\label{rem:bar-semi-infinite-scope}
The theorem constructs a lift from actual filtered coefficient data and compatible vanishing obstructions.
A comparison theorem does not supply those data from an isomorphism of first pages.
Example~\ref{ex:fl19-no-lift} proves that the proposed implication fails even for two-dimensional complexes.
Neither a state-level PBW basis nor conformal-weight preservation identifies the full collision sheaves with the zero-mode Chevalley complex.
The native map $f_0$, its higher corrections, and their compatibility with restrictions and differential operators remain additional constructions.
At critical level, a comparison must likewise specify its target and filtration independently of the undefined Sugawara stress tensor.
\end{remark}

\begin{remark}[Comparison with the string BRST complex]
\label{rem:string-brst-special-case}
To specialize~\eqref{eq:bar-semi-infinite} to a string BRST complex, one must identify its matter and ghost spaces, central action, charge term and differential with the chosen semi-infinite data.
The nilpotency condition of the string differential must hold on that carrier.
This identification is separate from the filtered lifting theorem.
\end{remark}

\subsection{Anomaly duality from complementarity}

\begin{corollary}[Anomaly duality for Kac--Moody pairs;
\ClaimStatusProvedHere]
\label{cor:anomaly-duality-km}
\index{anomaly cancellation!Kac--Moody duality}
\index{Feigin--Frenkel duality!anomaly interpretation}
Let $\widehat{\fg}_k$ be an affine Kac--Moody algebra at generic
level $k$, with Feigin--Frenkel dual
$\widehat{\fg}_{k'} = \widehat{\fg}_{-k-2h^\vee}$.  Then:
\begin{enumerate}[label=\textup{(\alph*)}]
\item \emph{Bar cohomology = semi-infinite cohomology.}
If the native bar--semi-infinite chain comparison is supplied,
the bar cohomology computations of Part~\textup{II} transport to
semi-infinite cohomology.  The generic-rank statement of
Theorem~\ref{thm:bar-cohomology-level-independence} does not compute a
Hilbert series or construct that comparison.  Under the chain-comparison
hypothesis, the dimension equality is
\[
\dim H^n(\barB^{\mathrm{ch}}(\widehat{\fg}_k))
= \dim H^{\infty/2+n}(\widehat{\fg}_k, V_k).
\]

\item \emph{Curvature = conformal anomaly.}
The genus-$1$ curvature coefficient
$\kappa(\widehat{\fg}_k) = (k+h^\vee)\dim\fg/(2h^\vee)$
\textup{(}Theorem~\textup{\ref{thm:genus-universality}(i))} equals
the semi-infinite anomaly coefficient.  The duality
$\kappa + \kappa' = 0$ of
Theorem~\textup{\ref{thm:genus-universality}(ii)} is the
anomaly cancellation for the Feigin--Frenkel pair: the
conformal anomalies of $\widehat{\fg}_k$ and
$\widehat{\fg}_{-k-2h^\vee}$ cancel.

\item \emph{Complementarity = anomaly duality.}
The quantum complementarity theorem
\textup{(}Theorem~\textup{\ref{thm:quantum-complementarity-main})}
applied to the Koszul pair
$(\widehat{\fg}_k, \widehat{\fg}_{-k-2h^\vee})$ gives:
\[
Q_g(\widehat{\fg}_k) \oplus Q_g(\widehat{\fg}_{-k-2h^\vee})
= H^*(\overline{\mathcal{M}}_g,\, Z(\widehat{\fg}_k)).
\]
Via the bar--semi-infinite identification, the left side is a
direct sum of semi-infinite cohomology spaces on
$\overline{\mathcal{M}}_g$: the genus-$g$ quantum corrections
of the WZW model at level~$k$ and its Feigin--Frenkel dual sum to
the full moduli space cohomology.  This is the mathematical
content of anomaly cancellation for Kac--Moody pairs at all genera.

\item \emph{Universal MC duality.}
The universal Maurer--Cartan classes satisfy
$\Theta_{\widehat{\fg}_k} + \Theta_{\widehat{\fg}_{-k-2h^\vee}}
= 0$ (Theorem~\ref{thm:explicit-theta}), since
$\kappa(\widehat{\fg}_k) + \kappa(\widehat{\fg}_{-k-2h^\vee}) = 0$
by Theorem~\ref{thm:genus-universality}(ii).
This is the chain-level mechanism underlying~(c):
the Verdier involution exchanges
$\Theta_{\widehat{\fg}_k}$ and $-\Theta_{\widehat{\fg}_{k'}}$,
producing the Lagrangian splitting of
Theorem~\ref{thm:quantum-complementarity-main}.
\end{enumerate}
\end{corollary}

\begin{proof}
Part~(a) is the cohomology-level consequence of
Theorem~\ref{thm:bar-semi-infinite-km}.  Part~(b): the
curvature $m_0 = \kappa \cdot \mathbf{1}$ in the bar complex
corresponds, under the quasi-isomorphism~$\Psi$, to the
semi-infinite anomaly (the failure of the naive CE differential
to square to zero, corrected by the semi-infinite charge).
The genus universality theorem identifies this with the
conformal anomaly coefficient.  The duality $\kappa + \kappa' = 0$
follows from the Feigin--Frenkel level shift and the formula
$\kappa = (k+h^\vee)\dim\fg/(2h^\vee)$: replacing $k$ by
$-k - 2h^\vee$ negates~$\kappa$.
Part~(c): combine the bar--semi-infinite identification with
Theorem~\ref{thm:quantum-complementarity-main}.  The genus-$g$
obstruction space $Q_g(\widehat{\fg}_k)$ is the associated
graded of $H^*(\barB^{(g)}(\widehat{\fg}_k))$, which by
Theorem~\ref{thm:bar-semi-infinite-km} is the genus-$g$
semi-infinite cohomology.  The complementarity sum equals
$H^*(\overline{\mathcal{M}}_g, Z(\widehat{\fg}_k))$ by
Theorem~\ref{thm:quantum-complementarity-main}.
\end{proof}

\begin{remark}[The Master Table, revisited]
\label{rem:master-table-brst}
\index{master table of invariants!BRST interpretation}
Via Corollary~\ref{cor:anomaly-duality-km}, every entry in the Master Table (Table~\ref{tab:master-invariants}) is a semi-infinite anomaly coefficient: $\kappa$ is the semi-infinite anomaly, $c + c' = 2\dim\fg$ the total Feigin--Frenkel anomaly, and $F_g$ the genus-$g$ WZW free energy.
\end{remark}

\subsection{Extension to $\mathcal{W}$-algebras via Drinfeld--Sokolov reduction}

\begin{theorem}[Bar complex = semi-infinite complex for \texorpdfstring{$\mathcal{W}$}{W}-algebras;
\ClaimStatusProvedHere]
\label{thm:bar-semi-infinite-w}
\index{semi-infinite cohomology!W-algebra identification}
\index{Drinfeld--Sokolov reduction!semi-infinite cohomology}
Let $\mathcal{W}^k(\fg, f_{\mathrm{prin}})$ be the principal
W-algebra at level $k \neq -h^\vee$.  There is a quasi-isomorphism:
\begin{equation}\label{eq:bar-semi-infinite-w}
\barB^{\mathrm{ch}}(\mathcal{W}^k(\fg))
\;\xrightarrow{\;\sim\;}\;
C^{\infty/2+\bullet}(\mathcal{W}^k(\fg),\, M_k)
\end{equation}
where $M_k$ is the vacuum module for $\mathcal{W}^k(\fg)$ and the
right side is the semi-infinite complex defined with respect to the
conformal weight grading.  In particular:
\[
H^*(\barB^{\mathrm{ch}}(\mathcal{W}^k(\fg)))
\;\cong\;
H^{\infty/2+*}(\mathcal{W}^k(\fg),\, M_k).
\]
\end{theorem}

\begin{proof}
The W-algebra is obtained from the Kac--Moody algebra by BRST
reduction (Definition~\ref{def:quantum-ds}):
$\mathcal{W}^k(\fg) = H^0(
\widehat{\fg}_k \otimes \mathcal{F}_{\mathrm{gh}},\,
Q_{\mathrm{DS}})$.
We compose Theorem~\ref{thm:bar-semi-infinite-km} with the exactness
of DS reduction on the relevant filtered complexes.

\emph{Step~1: Functoriality of the bar construction under DS.}
The bar functor sends $\widehat{\fg}_k$ to
$\barB^{\mathrm{ch}}(\widehat{\fg}_k)$, and the DS differential
$Q_{\mathrm{DS}}$ acts on the bar complex (since $Q_{\mathrm{DS}}$
is a derivation of the vertex algebra structure; the bar construction
preserves this action as it is functorial for $A_\infty$ morphisms).
By the proof of Theorem~\ref{thm:w-algebra-koszul-main}, Step~2,
the bar functor is exact on the homotopy category and:
\[
\barB^{\mathrm{ch}}(\mathcal{W}^k(\fg))
\;\simeq\;
H^*_{Q_{\mathrm{DS}}}\!\bigl(
\barB^{\mathrm{ch}}(\widehat{\fg}_k)
\otimes \barB(\mathcal{F}_{\mathrm{gh}})\bigr).
\]

\emph{Step~2: DS reduction of the semi-infinite complex.}
On the semi-infinite side, DS reduction of
$C^{\infty/2+\bullet}(\widehat{\fg}_k, V_k)$ yields
$C^{\infty/2+\bullet}(\mathcal{W}^k(\fg), M_k)$:
the semi-infinite complex of the W-algebra is the BRST
reduction of the semi-infinite complex of $\widehat{\fg}_k$.
This is the content of Feigin--Frenkel's construction of
W-algebra BRST cohomology~\cite{Feigin-Frenkel}: the
semi-infinite differential $Q_{\mathrm{si}}$ commutes with
$Q_{\mathrm{DS}}$ (both are BRST differentials for different
symmetries; their commutator vanishes by the standard double
complex argument).

\emph{Step~3: Comparison.}
The quasi-isomorphism $\Psi$ of
Theorem~\ref{thm:bar-semi-infinite-km} intertwines the
$Q_{\mathrm{DS}}$-action on both sides (since $\Psi$ is
constructed via the PBW spectral sequence, which respects the
adjoint action of $\widehat{\fg}_k$ on itself and hence the
DS differential).  Passing to $Q_{\mathrm{DS}}$-cohomology:
\[
H^*_{Q_{\mathrm{DS}}}(\Psi) \colon
H^*_{Q_{\mathrm{DS}}}(\barB^{\mathrm{ch}}(\widehat{\fg}_k)
\otimes \barB(\mathcal{F}_{\mathrm{gh}}))
\;\xrightarrow{\;\sim\;}\;
H^*_{Q_{\mathrm{DS}}}(
C^{\infty/2+\bullet}(\widehat{\fg}_k, V_k)
\otimes C^*(\mathcal{F}_{\mathrm{gh}}))
\]
is a quasi-isomorphism (a quasi-isomorphism of filtered complexes
remains a quasi-isomorphism after applying an exact functor on the
derived category, which $H^*_{Q_{\mathrm{DS}}}$ is for
$k \neq -h^\vee$ by Arakawa's vanishing theorem
\cite{Ara15}).  The left side is
$\barB^{\mathrm{ch}}(\mathcal{W}^k(\fg))$ (Step~1); the right
side is
$C^{\infty/2+\bullet}(\mathcal{W}^k(\fg), M_k)$ (Step~2).
\end{proof}

\begin{corollary}[Virasoro bar complex = semi-infinite complex;
\ClaimStatusProvedHere]
\label{cor:virasoro-semi-infinite}
\index{semi-infinite cohomology!Virasoro}
\index{Virasoro algebra!semi-infinite cohomology}
For the Virasoro algebra $\mathrm{Vir}_c$ at central charge
$c = 1 - 6(k+1)^2/(k+2)$ \textup{(}obtained as
$\mathcal{W}^k(\mathfrak{sl}_2)$\textup{)}:
\[
H^*(\barB^{\mathrm{ch}}(\mathrm{Vir}_c))
\;\cong\;
H^{\infty/2+*}(\mathrm{Vir}_c,\, M_c)
\]
where $M_c$ is the vacuum Virasoro module.  The Motzkin-number
bar cohomology dimensions of
Theorem~\textup{\ref{thm:virasoro-chiral-koszul}} are therefore
semi-infinite cohomology dimensions.
\end{corollary}

\begin{proof}
Apply Theorem~\ref{thm:bar-semi-infinite-w} with
$\fg = \mathfrak{sl}_2$, $f = f_{\mathrm{prin}}$.
\end{proof}

\begin{corollary}[Anomaly complementarity for \texorpdfstring{$\mathcal{W}$}{W}-algebra pairs;
\ClaimStatusProvedHere]
\label{cor:anomaly-duality-w}
\index{anomaly cancellation!W-algebra duality}
Let $\mathcal{W}^k(\fg)$ be the principal W-algebra, with
Feigin--Frenkel same-family partner
$\mathcal{W}^{k'}(\fg) = \mathcal{W}^{-k-2h^\vee}(\fg)$.
Then:
\begin{enumerate}[label=\textup{(\alph*)}]
\item The semi-infinite anomaly coefficients satisfy the same
level-independent complementarity law as the bar curvature
coefficients.  For principal finite-type $\mathcal{W}$-algebras,
\[
\kappa(\mathcal{W}^k(\fg)) + \kappa(\mathcal{W}^{k'}(\fg))
= \bigl(c(\mathcal{W}^k(\fg)) + c(\mathcal{W}^{k'}(\fg))\bigr)
\varrho(\fg),
\]
where $\varrho(\fg) = \sum_{i=1}^r \frac{1}{m_i+1}$ is the anomaly
ratio from Theorem~\ref{thm:wn-obstruction}.  Thus the sum is a
level-independent constant: $13$ for Virasoro and $250/3$ for
$\mathcal{W}_3$.

\item The quantum complementarity theorem
\textup{(}Theorem~\textup{\ref{thm:quantum-complementarity-main})}
applied to the W-algebra Koszul pair
\textup{(}Theorem~\textup{\ref{thm:w-algebra-koszul-main})} gives:
\[
Q_g(\mathcal{W}^k) \oplus Q_g(\mathcal{W}^{-k-2h^\vee})
= H^*(\overline{\mathcal{M}}_g,\, Z(\mathcal{W}^k)).
\]
Via the bar--semi-infinite identification, this is anomaly duality
for W-algebra pairs at all genera.

\item For $\fg = \mathfrak{sl}_2$
\textup{(}Virasoro\textup{)}: the Virasoro at
$c$ and its same-family shadow partner at $c' = 26 - c$ have
semi-infinite anomaly coefficients adding to~$13$, and their
genus-$g$ quantum corrections sum to
  $H^*(\overline{\mathcal{M}}_g, Z(\mathrm{Vir}_c))$.  Here
  $\mathrm{Vir}_{26-c}$ is the proved M/S-level shadow partner used in
  the semi-infinite calculation; the stronger H-level
  infinite-generator realization belongs to MC4
  (Chapter~\ref{chap:concordance}).
\end{enumerate}
\end{corollary}

\begin{proof}
Combine Theorem~\ref{thm:bar-semi-infinite-w} with
Theorem~\ref{thm:wn-obstruction} and
Theorem~\ref{thm:central-charge-complementarity}: the DS reduction
preserves the anomaly data, and the principal finite-type
$\mathcal{W}$-obstruction is $\kappa = c\,\varrho(\fg)$.
Part~(b) then follows from the quantum complementarity theorem.
Part~(c): specialize to $\fg = \mathfrak{sl}_2$ with $h^\vee = 2$,
for which $c + c' = 26$ and $\varrho(\mathfrak{sl}_2)=1/2$.
\end{proof}

\section{Holomorphic-topological field theories}

\subsection{Gaiotto's framework: from 4d to 2d}

\begin{definition}[Holomorphic-topological (HT) theory]
A holomorphic-topological field theory on a complex surface $\Sigma$ is holomorphic in one complex direction and topological in the other.  Fields are sections of $\mathcal{O}$-modules satisfying $\bar{\partial}_{\bar{z}} \phi = 0$ (holomorphic) and $Q_{\text{top}} \phi = 0$ (topologically BRST-closed).
\end{definition}

\begin{theorem}[HT theory from 4d \texorpdfstring{$\mathcal{N}=4$}{N=4} SYM; \ClaimStatusProvedElsewhere]
\label{thm:ht-from-n4-sym}
Starting with 4d $\mathcal{N}=4$ super Yang--Mills with gauge group $G$, one applies the holomorphic twist (Costello~\cite{costello-gaiotto}), localizes to $\bar{\partial}$-connections on a Riemann surface $\Sigma$, and obtains holomorphic BF theory (the dimensional reduction of holomorphic Chern--Simons to $\Sigma$).

The action is:
\[S_{\text{BF}} = \int_{\Sigma} \operatorname{Tr}\!\left(B \wedge \bar{\partial} \mathbf{A} + B \wedge \mathbf{A} \wedge \mathbf{A}\right)\]

where $\mathbf{A} = c + A + \cdots$ denotes the full BV field with $c \in \Omega^{0,0}(\Sigma, \mathfrak{g})$ (ghost) and $A \in \Omega^{0,1}(\Sigma, \mathfrak{g})$ (connection), and $B \in \Omega^{1,0}(\Sigma, \mathfrak{g})$ is the adjoint-valued Lagrange multiplier.
On a Riemann surface, $\Omega^{0,2}(\Sigma) = 0$, so the classical fields
$A \in \Omega^{0,1}$ satisfy $\bar{\partial}A = 0$ and $A \wedge A = 0$ identically.  The non-trivial
content of the action arises from the BV extension: the ghost component $c \in \Omega^{0,0}$ has
$\bar{\partial}c \in \Omega^{0,1}$, giving $B \wedge \bar{\partial}c \in \Omega^{1,1}$ (the volume form).
The full holomorphic Chern--Simons action $\int \Omega \wedge \operatorname{CS}(\mathcal{A})$ lives on a
complex $3$-fold with $(3,0)$-form $\Omega$; on $\Sigma$ only the holomorphic BF action above survives the dimensional reduction.
\end{theorem}

\begin{proof}[From Costello--Gaiotto]
\emph{Step~1: The Holomorphic Twist.}

Start with $\mathcal{N}=4$ SYM on $\mathbb{R}^4 \cong \mathbb{C}^2$. The field 
content includes the gauge field $A_\mu$, scalars $\Phi^I$ in the adjoint ($I=1,\ldots,6$), and fermions $\psi, \bar{\psi}$.  The twist redefines the Lorentz group action by mixing with R-symmetry:
\[\text{Spin}(4) \times \text{Spin}(6)_R \to \text{Spin}(4)_{\text{new}}\]

After twisting, some bosons become fermions (ghosts), some fermions become bosons (matter), and a nilpotent supercharge $Q$ becomes scalar.

\emph{Step~2: Localization.}

The twisted action has $Q^2 = 0$ and:
\[S_{\text{twisted}} = \{Q, \Lambda\} + S_0\]

By localization, path integral reduces to:
\[Z = \int_{\text{Q-fixed locus}} \mathcal{O}(fields)\]

The Q-fixed locus consists of holomorphic data:
\[\bar{\partial}_A = 0, \qquad \bar{\partial}_A \Phi = 0\]

These are the equations for a holomorphic bundle with holomorphic Higgs field (the Hitchin equations in the holomorphic sector).

\emph{Step~3: Volume Form.}

The holomorphic volume form $\Omega$ arises from the topological twist. On
$\mathbb{C}^2$ with holomorphic coordinates $(z,w)$:
\[\Omega = dz \wedge dw\]

On the deformed conifold $Y_\mu = \{u_1 w_2 - u_2 w_1 = \mu\} \subset \mathbb{C}^4$ (a Calabi--Yau 3-fold), $\Omega$ is the unique holomorphic $(3,0)$-form obtained via the Poincar\'{e} residue:
\[\Omega = \operatorname{Res}_{Y_\mu}\frac{du_1 \wedge du_2 \wedge dw_1 \wedge dw_2}{u_1 w_2 - u_2 w_1 - \mu}\]
\end{proof}

\subsection{Boundary conditions and chiral algebras}

\begin{theorem}[Boundary chiral algebra; \ClaimStatusProvedElsewhere{} \cite{CDG20,CG17}]
\label{thm:boundary-chiral-algebra-bv}
A boundary condition for holomorphic Chern--Simons theory supports a chiral
algebra $\mathcal{A}_{\text{bdy}}$ whose generators are local operators at the boundary, whose OPE comes from bulk-to-boundary correlation functions, and whose central charge is determined by the level of HCS theory.
\end{theorem}

\begin{example}[Kac--Moody from HCS]
Holomorphic Chern--Simons with gauge group $G$ at level $k$ produces:
\[\mathcal{A}_{\text{bdy}} = \widehat{\mathfrak{g}}_k\]

the affine Kac--Moody algebra at level $k$.

The currents:
\[J^a(z) = \text{Tr}(T^a A(z))\]

satisfy:
\[J^a(z) J^b(w) \sim \frac{k \delta^{ab}}{(z-w)^2} + 
\frac{f^{abc} J^c(w)}{z-w}\]

This OPE is computed via the holomorphic Chern--Simons path integral with
boundary insertions.
\end{example}

\subsection{The holomorphic-topological boundary condition}

\begin{definition}[HT boundary condition]
For holomorphic Chern--Simons on a surface $\Sigma$ with boundary $\partial \Sigma$, 
the \emph{holomorphic-topological boundary condition} requires:
\begin{enumerate}
\item Fields extend holomorphically to a compactification $\bar{\Sigma}$
\item At the boundary divisor $D = \bar{\Sigma} \setminus \Sigma$, fields have 
a simple zero: $A \in \Omega^{0,*}(\bar{\Sigma}, \mathcal{O}_{\bar{\Sigma}}(-D))$
\item The volume form $\Omega$ has compatible pole: $\Omega \in K_{\bar{\Sigma}}(kD)$ 
for cubic interaction $(k=3)$
\end{enumerate}
\end{definition}

\begin{conjecture}[Bar-cobar from HT boundary; \ClaimStatusHeuristic]
\label{conj:bar-cobar-ht-boundary}
The holomorphic-topological boundary condition realizes bar-cobar duality:
\begin{itemize}
\item \emph{Bar side.} Fields with prescribed asymptotics near $D$ (logarithmic
forms)
\item \emph{Cobar side.} Distributional fields on $\Sigma$ (delta functions at $D$)
\item \emph{Duality.} Perfect pairing via residue-distribution integral
\end{itemize}
\end{conjecture}

\begin{evidence}[Evidence: Geometric realization]
The key is the volume form behavior. If $\Omega$ has a pole of order $k$ along 
$D$:
\[\Omega \sim \frac{dz \wedge dw}{(z-w)^k}\]

then the action term:
\[\int_{\Sigma} \Omega \cdot A^k\]

is finite if and only if $A$ vanishes to order 1 along $D$.

This pole-zero compatibility is exactly the bar-cobar relationship: the logarithmic form $\eta = d\log(z-w)$ (bar) has a logarithmic singularity whose residue is the distribution $\delta(z-w)$ (cobar).

The holomorphic-topological boundary condition enforces this duality at the
geometric level.
\end{evidence}

\begin{remark}[Scope]
This conjecture is labeled \ClaimStatusHeuristic{} because it proposes a geometric
realization of bar-cobar duality via holomorphic-topological boundary conditions,
which requires input from extended field theory (boundary conditions for HT theories)
beyond the algebraic framework of this monograph.

\emph{Homotopy template} (Convention~\ref{conv:hms-levels}): Type~VII (physics-dictionary).
The bar and cobar complexes should arise as the two boundary conditions (logarithmic forms vs.\ distributional fields) of a holomorphic-topological theory on~$\Sigma$, with the bar-cobar pairing realized by the residue-distribution integral across the boundary divisor~$D$.
\end{remark}

\section{$\mathcal{W}$-algebras from Higgs branches}

\subsection{\texorpdfstring{4d gauge theory $\to$ 2d $\mathcal{W}$-algebra}{4d gauge theory -> 2d $\mathcal{W}$-algebra}}

\begin{conjecture}[Costello--Gaiotto AGT; \ClaimStatusConjectured]
\label{conj:costello-gaiotto-agt}
Starting with 4d $\mathcal{N}=2$ gauge theory with gauge group $G$, compactify on $\Sigma_g$, apply the topological twist, and take the infrared limit.  The result is a 2d CFT with W-algebra symmetry $\mathcal{W}(G)$.

For $G = SU(N)$, this gives the $\mathcal{W}_N$ algebra.
\end{conjecture}

\begin{evidence}[Evidence: Via bar-cobar]
\emph{Step~1: Higgs Moduli Space.}

The 4d theory has Higgs branch moduli space:
\[\mathcal{M}_{\text{Higgs}} = \text{Higgs}(\Sigma_g, G)\]

consisting of $G$-Higgs bundles on $\Sigma_g$.

\emph{Step~2: Chiral Algebra.}

The Higgs moduli space supports a chiral algebra via:
\[\mathcal{A} = \text{LocalObs}(\mathcal{M}_{\text{Higgs}})\]

Local operators on the moduli space form a factorization algebra, which extends 
to a chiral algebra.

\emph{Step~3: W-Algebra Identification.}

For $G = SU(N)$ and $\Sigma_g = \mathbb{C}$, the local operators --- the stress tensor $T(z)$ from diffeomorphisms and higher spin currents $W^{(s)}(z)$ for $s = 2, \ldots, N$ --- generate the $\mathcal{W}_N$ algebra.

\emph{Step~4: Bar-Cobar Realization.}

The bar complex:
\[\bar{B}^{\text{ch}}(\mathcal{W}_N) = \Omega^*(\overline{C}_n(\Sigma_g), 
\mathcal{W}_N^{\boxtimes n})\]

computes the boundary W-algebra correlators on the 2d side.  The
conjectural gauge-theory statement is that these same correlators arise
from partition functions of the 4d theory:
\[\langle W^{(s_1)}(z_1) \cdots W^{(s_n)}(z_n) \rangle =
Z_{4d}[\Sigma_g; z_1, \ldots, z_n]\]
\end{evidence}

\begin{remark}[Scope]
\ClaimStatusConjectured{} because the 4d gauge theory origin of W-algebras (Costello--Gaiotto) requires the holomorphic twist and infrared localization, which are physics inputs beyond the scope of this monograph.  The bar complex computing W-algebra conformal blocks follows from Part~I; the live task is the gauge-theoretic comparison.

\emph{Homotopy template} (Convention~\ref{conv:hms-levels}): Type~VII (physics-dictionary).  Contributing to Conjecture~\ref{conj:master-bv-brst}.
\end{remark}

\subsection{Quantum corrections and central charge}

\begin{remark}[Quantum vs.\ classical]
The 4d $\to$ 2d reduction involves two types of quantum corrections: \emph{loop corrections} from integrating out massive modes (captured by $m_3, m_4, \ldots$ in the $A_\infty$ structure), and \emph{instanton corrections} from non-perturbative effects (not visible in the bar complex, requiring full QFT).
\end{remark}

\begin{theorem}[Central charge from 4d; \ClaimStatusProvedElsewhere]
\label{thm:central-charge-from-4d}
The Sugawara central charge of the affine Kac--Moody algebra $\widehat{\mathfrak{g}}_k$ is:
\[c_{\mathrm{Sug}} = \frac{k \dim \mathfrak{g}}{k + h^\vee}\]

The W-algebra $\mathcal{W}^k(\mathfrak{g}, f)$ obtained by DS reduction has a different central charge, determined by $c_{\mathrm{Sug}}$ minus the ghost contribution (which depends on the nilpotent orbit of $f$), where $k$ is the level and $h^\vee$ the dual Coxeter number.  This matches the Feigin--Frenkel formula~\cite{Feigin-Frenkel}.
\end{theorem}

\section{Quantum observables and BV integration}

\subsection{BV path integral}

\begin{definition}[BV partition function]
The BV partition function is:
\[Z_{\text{BV}} = \int_{\mathcal{L}} [D\phi] \, e^{S[\phi]/\hbar}\]

where $\mathcal{L}$ is a Lagrangian submanifold (gauge fixing), $S[\phi]$ satisfies the quantum master equation $\Delta_{\text{BV}} e^{S/\hbar} = 0$, and the integration uses the BV measure (Berezinian).
\end{definition}

\begin{conjecture}[BV integration = bar-cobar pairing; \ClaimStatusHeuristic]
\label{conj:bv-integration-bar-cobar}
The BV path integral is realized by the bar-cobar pairing:
\[Z_{\text{BV}} = \langle \bar{B}^{\text{ch}}(\mathcal{A}), 
\Omega^{\text{ch}}(\mathcal{C}) \rangle\]

Explicitly:
\[Z = \int_{\overline{C}_n(X)} \omega_{\text{bar}} \wedge \iota^* K_{\text{cobar}}\]
\end{conjecture}

\begin{evidence}
\emph{Step~1: Gauge Fixing as Lagrangian.}

Choose gauge fixing Lagrangian $\mathcal{L}$ corresponding to:
\[\mathcal{L} = \{(\phi, \phi^+) : \phi^+ = F(\phi)\}\]

for some gauge fermion $F$.

In geometric terms, this corresponds to choosing a regularization prescription 
for the configuration space integrals.

\emph{Step~2: BV Measure.}

The BV measure on $\mathcal{L}$ is:
\[\mu_{\text{BV}} = \text{Ber}(\mathcal{L}) \cdot d\phi\]

Geometrically, this is the measure on configuration space:
\[\mu_{\text{geom}} = \prod_{i<j} |z_i - z_j|^2 \prod_{i=1}^n d^2z_i\]

with appropriate gauge fixing factors.

\emph{Step~3: Action and Pairing.}

The action:
\[e^{S/\hbar} = \prod_{\text{interactions}} e^{V_k(z_1, \ldots, z_k)/\hbar}\]

corresponds to the cobar element:
\[K_{\text{cobar}} = \sum_n K_n(z_1, \ldots, z_n)\]

The pairing:
\[\int_{\overline{C}_n(X)} \omega_{\text{bar}} \wedge K_{\text{cobar}}\]

is exactly the BV path integral with gauge fixing determined by the choice of
regularization.
\end{evidence}

\begin{remark}[Scope]
This conjecture is labeled \ClaimStatusHeuristic{} because the identification of the
BV path integral with the bar-cobar pairing requires a rigorous measure-theoretic
framework for configuration space integrals with BV gauge-fixing data, which goes
beyond the algebraic constructions of Part~I.

\emph{Homotopy template} (Convention~\ref{conv:hms-levels}): Type~VII (physics-dictionary).
The BV partition function $Z_{\mathrm{BV}}$ should equal the bar-cobar pairing $\langle \bar{B}^{\mathrm{ch}}(\cA), \Omega^{\mathrm{ch}}(\mathcal{C}) \rangle$ as a configuration space integral, with gauge-fixing Lagrangian corresponding to the choice of regularization of the FM compactification.
\end{remark}

\subsection{Observables and correlation functions}

\begin{theorem}[Observables = cohomology; \ClaimStatusProvedElsewhere{} \cite{CG17}]
\label{thm:observables-brst-cohomology}
Physical observables are:
\[\mathcal{O}_{\text{phys}} = H^0(Q_{\text{BRST}}) = 
\ker(Q_{\text{BRST}})/\text{im}(Q_{\text{BRST}})\]

In the bar-cobar framework:
\[\mathcal{O}_{\text{phys}} \cong H^0(\bar{B}^{\text{ch}}(\mathcal{A}))\]
\end{theorem}

\begin{example}[Correlation functions]
An $n$-point correlation function:
\[\langle \mathcal{O}_1(z_1) \cdots \mathcal{O}_n(z_n) \rangle = 
\int_{\overline{C}_n(X)} [\mathcal{O}_1 \otimes \cdots \otimes \mathcal{O}_n] 
\cdot e^{S_{\text{int}}}\]

is computed by representing each $\mathcal{O}_i$ as a cocycle in $\bar{B}^{\text{ch}}(\mathcal{A})$ and applying the bar-cobar pairing with $e^{S_{\text{int}}}$ --- Gaiotto's prescription for holomorphic-topological observables.
\end{example}

\section{Dictionary}

\begin{remark}[Dictionary]
The BV-BRST formalism connects:

\begin{center}
\begin{tabular}{|l|l|l|}
\hline
\textbf{Physics} & \textbf{Math (Bar-Cobar)} & \textbf{Geometry} \\
\hline
BV complex & Bar complex & Log forms on $\overline{C}_n$ \\
BV bracket & Configuration space symplectic & Poisson structure \\
Master equation & $\dzero^2 = 0$ & Stokes / boundary residues \\
BV Laplacian & Cobar differential & Delta functions \\
Quantum master eq & Bar-cobar duality & Verdier duality \\
BRST operator & Residue extraction & OPE singularities \\
Gauge fixing & Lagrangian choice & Regularization \\
Observables & Cohomology & Physical states \\
Path integral & Bar-cobar pairing & Configuration integrals \\
4d $\to$ 2d reduction & Factorization & Dimensional analysis \\
W-algebra & Boundary chiral algebra & Higgs moduli \\
\hline
\end{tabular}
\end{center}
\end{remark}

\begin{remark}[Gaiotto's contribution]
Gaiotto's contribution was recognizing a physical pattern: holomorphic-topological theories naturally produce chiral algebras, boundary conditions for these theories are modules over the chiral algebra, the open-closed correspondence should be modeled by bar-cobar duality, and suitable 4d gauge theory reductions should produce W-algebras via this mechanism.  The geometric bar-cobar construction supplies the theorematic boundary algebraic package on which these physical comparison statements are built.
\end{remark}


\section{The BV algebra structure}
\label{sec:complete-bv-structure}

\subsection{BV algebra definition}

\begin{definition}[BV algebra]
\label{def:bv-algebra-complete}
\index{BV algebra|textbf}
A \emph{Batalin--Vilkovisky algebra} is a graded commutative algebra $(A, \cdot)$ equipped with:
\begin{enumerate}
\item \emph{BV bracket.} $\{\cdot, \cdot\}: A \otimes A \to A$ of degree $+1$, satisfying:
\begin{itemize}
\item Graded skew-symmetry: $\{a, b\} = -(-1)^{(|a|+1)(|b|+1)} \{b, a\}$
\item Graded Leibniz: $\{a, bc\} = \{a, b\}c + (-1)^{(|a|+1)|b|} b\{a, c\}$
\item Graded Jacobi: $\{a, \{b, c\}\} = \{\{a, b\}, c\} + (-1)^{(|a|+1)(|b|+1)} \{b, \{a, c\}\}$
\end{itemize}

\item \emph{BV Laplacian.} $\Delta: A \to A$ of degree $+1$, satisfying:
\begin{itemize}
\item Nilpotency: $\Delta^2 = 0$
\item Second-order: $\Delta(ab) = \Delta(a)b + (-1)^{|a|}a\Delta(b) + (-1)^{|a|}\{a, b\}$
\end{itemize}

\item \emph{Compatibility.} $\{a, b\} = (-1)^{|a|} [\Delta(ab) - \Delta(a)b - (-1)^{|a|}a\Delta(b)]$
\end{enumerate}
\end{definition}

\subsection{BV structure from configuration spaces}

\begin{theorem}[Configuration space BV structure; \ClaimStatusProvedHere]
\label{thm:config-space-bv}
The bar complex $\bar{B}^{\mathrm{ch}}(\mathcal{A})$ carries a natural BV algebra structure where:

\emph{(1) Algebra structure.} Wedge product of logarithmic forms

\emph{(2) BV bracket.} Derived from symplectic structure on $T^*C_n(X)$

\emph{(3) BV Laplacian.} Integration against diagonal (cobar differential)
\[\Delta(\omega) = \sum_{i < j} \int_{\Delta_{ij}} \omega \cdot \delta(z_i - z_j)\]
\end{theorem}

\begin{proof}
We verify the three BV algebra axioms (Definition~\ref{def:bv-algebra-complete}).

\emph{(1) Commutative algebra structure.}
The bar complex $\bar{B}^n(\mathcal{A}) = \Gamma(\overline{C}_{n+1}(X), \mathcal{A}^{\boxtimes(n+1)} \otimes \Omega^*_{\log})$ carries a natural commutative product: the wedge product of logarithmic forms.  For $\omega \in \bar{B}^p$ and $\eta \in \bar{B}^q$, define:
\[
(\omega \cdot \eta)(z_0, \ldots, z_{p+q+1}) = \omega(z_0, \ldots, z_p) \wedge \eta(z_{p+1}, \ldots, z_{p+q+1})
\]
extended to $\overline{C}_{p+q+2}(X)$ via the inclusion $\overline{C}_{p+1}(X) \times \overline{C}_{q+1}(X) \hookrightarrow \overline{C}_{p+q+2}(X)$.  This is graded commutative: $\omega \cdot \eta = (-1)^{|\omega||\eta|}\, \eta \cdot \omega$ by the sign rule for wedge products of forms.

\emph{(2) BV Laplacian.}
Define $\Delta: \bar{B}^n \to \bar{B}^{n-1}$ by contraction with the diagonal:
\[
\Delta(\omega) = \sum_{i < j} \int_{\Delta_{ij}} \omega \cdot \delta(z_i - z_j) = \sum_{i < j} \operatorname{Res}_{z_i = z_j}[\omega]
\]
This has degree $+1$ (reduces bar degree by 1, increases cohomological degree by 1 via the residue).

\emph{Nilpotency} ($\Delta^2 = 0$): Applying $\Delta$ twice gives a double residue $\operatorname{Res}_{z_i = z_j} \circ \operatorname{Res}_{z_k = z_l}$.  For $\{i,j\} \cap \{k,l\} = \emptyset$, the two residues commute and $\Delta^2 = 0$ is immediate.  For $\{i,j\} \cap \{k,l\} \neq \emptyset$, say $j = k$, the iterated residue $\operatorname{Res}_{z_i = z_j}\operatorname{Res}_{z_j = z_l}$ picks out the coefficient of $(z_i - z_j)^{-1}(z_j - z_l)^{-1}$ in $\omega$, but the Arnold relation $\eta_{ij} \wedge \eta_{jl} + \eta_{jl} \wedge \eta_{li} + \eta_{li} \wedge \eta_{ij} = 0$ ensures these contributions cancel in pairs: the terms from triples $(i,j,l)$ with different orderings of the iterated residue sum to zero.

\emph{Second-order property}: We verify $\Delta(\omega \cdot \eta) = \Delta(\omega) \cdot \eta + (-1)^{|\omega|}\omega \cdot \Delta(\eta) + (-1)^{|\omega|}\{\omega, \eta\}$ by computing the residue of a product.  When $\operatorname{Res}_{z_i = z_j}$ acts on $\omega(z_0, \ldots, z_p) \wedge \eta(z_{p+1}, \ldots, z_{p+q+1})$, the residue localizes to three cases: (a) both $i, j \leq p$, contributing $\Delta(\omega) \cdot \eta$; (b) both $i, j > p$, contributing $\omega \cdot \Delta(\eta)$; (c) $i \leq p < j$ or $j \leq p < i$, contributing the bracket $\{\omega, \eta\}$.  This decomposition is the BV second-order property.

\emph{(3) BV bracket.}
The bracket is defined as the failure of $\Delta$ to be a derivation:
\[
\{\omega, \eta\} = (-1)^{|\omega|}[\Delta(\omega \cdot \eta) - \Delta(\omega) \cdot \eta - (-1)^{|\omega|}\omega \cdot \Delta(\eta)]
\]
Geometrically, $\{\omega, \eta\}$ extracts the residue across the interface between the two factors — this is the contraction pairing between $\omega$ and $\eta$ mediated by the symplectic form on $T^*C_n(X)$.  The Jacobi identity for $\{-,-\}$ follows from $\Delta^2 = 0$ by the standard BV$\to$Lie derivation (the ``7-term identity'').

Thus $(\bar{B}^{\mathrm{ch}}(\mathcal{A}), \cdot, \Delta, \{-,-\})$ is a BV algebra.
\end{proof}

\subsection{Quantum master equation}

\begin{conjecture}[Quantum master equation; \ClaimStatusHeuristic]
\label{conj:quantum-master-complete}
The \emph{quantum master equation}
\[\hbar \Delta S + \frac{1}{2}\{S, S\} = 0\]
or equivalently $\Delta e^{S/\hbar} = 0$
is solved by the bar-cobar pairing.
\end{conjecture}

\begin{remark}[Scope]
This conjecture is labeled \ClaimStatusHeuristic{} because the solution $S$ of
the QME must be constructed from the bar-cobar pairing data, which requires
controlling the $\hbar$-expansion of configuration space integrals and verifying
the second-order BV identity at each order.

\emph{Homotopy template} (Convention~\ref{conv:hms-levels}): Type~VII (physics-dictionary).
The bar-cobar pairing element $S = \langle \bar{B}^{\mathrm{ch}}(\cA), \Omega^{\mathrm{ch}}(\cA^!) \rangle$ should satisfy the quantum master equation $\hbar \Delta S + \tfrac{1}{2}\{S, S\} = 0$ in the BV algebra of Theorem~\ref{thm:config-space-bv}.
\end{remark}

\begin{conjecture}[BV quantization = bar-cobar duality; \ClaimStatusHeuristic]
\label{conj:bv-equals-bar-cobar}
The BV quantization of a chiral algebra $\mathcal{A}$ is equivalent to computing the
bar-cobar homology:
\[H^*_{\text{BV}}(\mathcal{A}) \cong H^*(\bar{B}(\mathcal{A}), \Omega(\mathcal{A}^!))\]
where $\mathcal{A}^!$ is the Koszul dual.

For affine Kac--Moody algebras, Theorem~\ref{thm:bar-semi-infinite-km} gives a chain comparison under its supplied coefficient, target, filtration, symbol-map, and compatible-correction hypotheses.
A BV-homology identification additionally requires the actual BV complex and compatibility of that map with the required BV operations and pairings.
The bar-cobar quasi-isomorphism of Theorem~\ref{thm:bar-cobar-isomorphism-main} does not supply these separate maps.
For general chiral algebras, their construction remains part of the conjecture.
\end{conjecture}

\begin{remark}[Scope]
For Kac--Moody algebras, the bar/BRST comparison retains all source, target, filtration, symbol-map, and correction hypotheses of Theorem~\ref{thm:bar-semi-infinite-km}.
Its BV compatibility and pairings require additional maps even when the bar-cobar quasi-isomorphism of Theorem~\ref{thm:bar-cobar-isomorphism-main} is available.
The conjecture includes those constructions and their extension to general chiral algebras.

\emph{Homotopy template} (Convention~\ref{conv:hms-levels}): Type~II (equivalence of homotopy categories).
The BV cohomology $H^*_{\mathrm{BV}}(\cA)$ and the bar-cobar cohomology $H^*(\bar{B}(\cA), \Omega(\cA^!))$ should be isomorphic as graded vector spaces for every Koszul chiral algebra, induced by a quasi-isomorphism of the underlying cochain complexes.
\end{remark}

\subsection{Summary: BV as functor}

\begin{theorem}[BV functor; \ClaimStatusProvedHere]
\label{thm:bv-functor}
\index{BV functor}
The BV quantization defines a functor:
\[\mathrm{BV}: \mathrm{ChirAlg}_X \longrightarrow \mathrm{BV\text{-}Alg}\]

preserving tensor products up to quasi-isomorphism (lax monoidal structure) and Verdier duality for Koszul $\mathcal{A}$: $\mathbb{D}(\bar{B}(\mathcal{A})) \cong \Omega(\mathcal{A}^!)$.
\end{theorem}

\begin{proof}
\emph{Objects.}  For a chiral algebra $\mathcal{A}$ on $X$, define $\mathrm{BV}(\mathcal{A}) = (\bar{B}^{\mathrm{ch}}(\mathcal{A}), \cdot, \Delta, \{-,-\})$, the BV algebra structure of Theorem~\ref{thm:config-space-bv}.  This assignment is well-defined: the bar complex, wedge product, and BV Laplacian depend only on the chiral algebra structure of~$\mathcal{A}$.

\emph{Morphisms.}  A morphism $f: \mathcal{A} \to \mathcal{B}$ of chiral algebras induces $\bar{B}(f): \bar{B}^{\mathrm{ch}}(\mathcal{A}) \to \bar{B}^{\mathrm{ch}}(\mathcal{B})$ via $f^{\boxtimes(n+1)}$ on each bar degree.  This map preserves the wedge product (it acts on coefficients, not on forms) and commutes with the BV Laplacian (which is defined by residues, and $f$ preserves OPE data).  Therefore $\bar{B}(f)$ is a morphism of BV algebras.  Functoriality ($\bar{B}(g \circ f) = \bar{B}(g) \circ \bar{B}(f)$ and $\bar{B}(\mathrm{id}) = \mathrm{id}$) is inherited from the functoriality of $\bar{B}^{\mathrm{ch}}$.

\emph{Monoidal structure.}  For chiral algebras $\mathcal{A}_1, \mathcal{A}_2$, the Eilenberg--Zilber theorem provides a quasi-isomorphism
\[
\bar{B}^{\mathrm{ch}}(\mathcal{A}_1 \otimes \mathcal{A}_2) \xrightarrow{\sim} \bar{B}^{\mathrm{ch}}(\mathcal{A}_1) \otimes \bar{B}^{\mathrm{ch}}(\mathcal{A}_2)
\]
via the shuffle map.  This is \emph{not} an isomorphism at the chain level: the bar construction involves iterated tensor powers over a single configuration space $\overline{C}_n(X)$, which does not factor as a product.  The shuffle quasi-isomorphism is compatible with the BV structures
in the following sense.  On the target, the BV Laplacian satisfies
$\Delta_{\otimes} = \Delta_1 \otimes \mathrm{id}
+ \mathrm{id} \otimes \Delta_2$
(the Laplacian distributes over the tensor product because it acts
by contraction against the diagonal, which factors).
The BV bracket
$\{a_1 \otimes a_2,\, b_1 \otimes b_2\}
= \{a_1, b_1\} \otimes a_2 b_2
\pm a_1 b_1 \otimes \{a_2, b_2\}$
satisfies the Leibniz rule by definition of the Gerstenhaber bracket
on a tensor product of BV algebras.  The shuffle map intertwines
these structures because it is a coalgebra map (it respects the
deconcatenation coproduct, which determines the Laplacian) and an
algebra map up to homotopy (the Eilenberg--Zilber theorem).
Explicitly, the failure of strict compatibility is controlled by
higher homotopies in the $E_\infty$-structure on the bar complex
(Berger--Fresse~\cite[\S11]{BerFres04}), which do not affect the
lax monoidal structure at the derived level.  The functor
$\mathrm{BV}$ is therefore lax monoidal.

\emph{Verdier duality.}  Verdier duality $\mathbb{D}$ on $\overline{C}_n(X)$ exchanges logarithmic forms (bar data) with distributional sections (cobar data).  Specifically, $\mathbb{D}(\bar{B}^n(\mathcal{A})) = \mathrm{Dist}(C_n(X), (\mathcal{A}^*)^{\boxtimes n}) \cong \Omega^n(\mathcal{A}^!)$ (the cobar complex of the Koszul dual coalgebra), where the isomorphism uses the residue pairing between forms and distributions.  Under this identification, the bar differential becomes the cobar differential of $\mathcal{A}^!$, and the BV Laplacian becomes the cobar product, establishing $\mathbb{D}(\bar{B}(\mathcal{A})) \cong \Omega(\mathcal{A}^!)$ as BV algebras.  This is the BV-algebraic manifestation of the bar-cobar Verdier duality proved in Theorem~\ref{thm:verdier-bar-cobar}.
\end{proof}

\begin{remark}[Shifted symplectic structure from BV]\label{rem:bv-shifted-symplectic}
\index{shifted symplectic!from BV}
The BV bracket $\{-,-\}$ (degree~$+1$) on $\barB^{\mathrm{ch}}(\cA)$
is a $(-1)$-shifted Poisson structure; the Verdier duality
identification $\mathbb{D}(\barB(\cA)) \cong \Omega(\cA^!)$
provides non-degeneracy.  By the equivalence between non-degenerate
shifted Poisson and shifted symplectic structures
\cite{CPTVV17}, this makes the formal moduli problem
$\mathrm{Def}_g(\cA)$ a $(-1)$-shifted symplectic derived stack.
The deformation-obstruction decomposition
$Q_g(\cA) \oplus Q_g(\cA^!)$ is a Lagrangian polarization.
This is the content of
Theorem~\ref{thm:shifted-symplectic-complementarity}.
\end{remark}
